\documentclass[11pt]{article}

% ---------- Packages ----------
\usepackage[utf8]{inputenc}
\usepackage[T1]{fontenc}
\usepackage[margin=1in]{geometry}
\usepackage{amsmath,amssymb,amsthm,mathtools}
\usepackage{booktabs}
\usepackage{enumitem}
\usepackage{xcolor}
\usepackage{titlesec}
\usepackage{longtable}
\usepackage{array}
\usepackage{tocloft}
\usepackage[colorlinks=true,linkcolor=Blue,citecolor=OliveGreen,urlcolor=BrickRed,linktoc=all]{hyperref}
\usepackage[capitalise,noabbrev]{cleveref}

% ---------- Colours ----------
\definecolor{Blue}{rgb}{0.0,0.2,0.55}
\definecolor{OliveGreen}{rgb}{0.15,0.4,0.15}
\definecolor{BrickRed}{rgb}{0.55,0.15,0.1}
\definecolor{proved}{rgb}{0.0,0.35,0.0}
\definecolor{claimed}{rgb}{0.55,0.35,0.0}
\definecolor{negative}{rgb}{0.55,0.1,0.1}

% ---------- Theorem environments ----------
\theoremstyle{plain}
\newtheorem{theorem}{Theorem}[section]
\newtheorem{proposition}[theorem]{Proposition}
\newtheorem{lemma}[theorem]{Lemma}
\newtheorem{corollary}[theorem]{Corollary}
\newtheorem{claimstmt}[theorem]{Claim}
\newtheorem{conjecture}[theorem]{Conjecture}
\theoremstyle{definition}
\newtheorem{definition}[theorem]{Definition}
\newtheorem{example}[theorem]{Example}
\theoremstyle{remark}
\newtheorem{remark}[theorem]{Remark}

% ---------- Status tags ----------
\newcommand{\LEAN}{\textcolor{proved}{\textsc{[Lean]}}}
\newcommand{\CLAIMED}{\textcolor{claimed}{\textsc{[Claimed]}}}
\newcommand{\WITHDRAWN}{\textcolor{negative}{\textsc{[Withdrawn]}}}
\newcommand{\NEGATIVE}{\textcolor{negative}{\textsc{[Negative result]}}}
\newcommand{\MEASURED}{\textcolor{claimed}{\textsc{[Measured]}}}
\newcommand{\file}[1]{\texttt{#1}}
\newcommand{\lean}[1]{\texttt{#1}}
\renewcommand{\_}{\textunderscore\allowbreak}
\sloppy
\emergencystretch=2.5em

% ---------- Misc ----------
\renewcommand{\thefootnote}{\arabic{footnote}}
\setlength{\cftbeforesecskip}{4pt}
\titlespacing*{\section}{0pt}{2.2ex plus 1ex minus .2ex}{1.4ex plus .2ex}
\titlespacing*{\subsection}{0pt}{1.8ex plus 1ex minus .2ex}{1ex plus .2ex}

\newcommand{\Z}{\mathbb{Z}}
\newcommand{\Q}{\mathbb{Q}}
\newcommand{\N}{\mathbb{N}}
\newcommand{\R}{\mathbb{R}}
\newcommand{\Fq}{\mathbb{F}_q}
\newcommand{\Br}{\operatorname{Br}}
\newcommand{\inv}{\operatorname{inv}}
\DeclareMathOperator{\ord}{ord}

\title{\textbf{The Riley Betts Erd\H{o}s--Straus Programme:\\[2pt]
A Lean~4 Formalisation and Analytic-Geometric Investigation\\[2pt]
\large Project Report on the Public Archive \texttt{erdos-straus-main}\\[2pt]
\normalsize Revision v0.3}}

\author{Prepared for Riley Betts Ltd \\ \small Compiled from the public git archive of the
Erd\H{o}s--Straus programme \\ \small (\texttt{erdos-straus-main}, 21 August 2026)}
\date{22 August 2026}

\begin{document}
\maketitle

\begin{center}
\fbox{\parbox{0.92\textwidth}{\small\textbf{What is new in v0.3 (22 August 2026).}
This revision supersedes v0.2. The E\_power theorem of v0.2
(\(S_A(x,2x)\ll x^{1-\delta}\) via one-stage hub-conditioned Suen plus
an \(O(M_T)\) transfer) is \WITHDRAWN{} as a theorem. A two-stage
replacement at a prime-power or largest-prime residual hub failed both
revival numbers through \(A=32\). Current record:
\texttt{erdos-straus-E-power.md}, \texttt{erdos-straus-E-power-decision.md}.
The combinatorial core in \texttt{EPower.lean} stands.
Vaughan (1970) remains the published comparison.}}
\end{center}
\vspace{1em}

\begin{abstract}
This report documents the full contents of \file{erdos-straus-main}, the public git archive of
an extensive computer-assisted research programme directed at the Erd\H{o}s--Straus conjecture.
The programme does \emph{not} claim a proof. Its concrete output is threefold. First, a
substantial Lean~4 formalisation (thirteen source files, several thousand lines, axiom-audited)
that reduces the conjecture to a single unproved analytic interface, proves a range of
elementary and Brauer--Manin-theoretic facts about the associated affine surfaces, and certifies
a finite verified range. Second, a sequence of analytic-number-theoretic partial results and
honestly labelled negative findings on a covering-density reformulation of the conjecture,
including a later-withdrawn one-stage covering-box claim (see the dated
notice below), and
a capstone diagnosis of why four independent completion attempts share a single obstruction.
Third, a geometric investigation of the Bright--Loughran Brauer--Manin approach, which is shown
to reduce, together with every other geometric mechanism considered, to the same analytic core.
The archive also triages a sibling research effort and an external claimed proof, keeping only
what survives scrutiny as library material. We describe the mathematical objectives, present the
formal architecture in detail, discuss each analytic and geometric approach investigated
together with its results (positive, partial, or negative), catalogue the accompanying numerical
experiments, and give a complete inventory of every file in the archive together with the full
bibliography the programme draws on.
\end{abstract}

\tableofcontents


\section{Introduction}
\label{sec:intro}

\subsection{The Erd\H{o}s--Straus conjecture}

The Erd\H{o}s--Straus conjecture (ES), posed in 1948 and still open, asserts that for every
integer $n\ge 2$ there exist positive integers $x,y,z$ with
\[
\frac{4}{n} = \frac{1}{x}+\frac{1}{y}+\frac{1}{z}.
\]
Clearing denominators gives the equivalent Diophantine statement
\[
\mathrm{IsES}(n,x,y,z) :\iff\; x,y,z>0 \ \text{and}\ 4xyz = n(xy+yz+zx).
\]
Classical work of Mordell, Schinzel, Aigner, Rosati, Yamamoto and others reduces the conjecture,
via explicit polynomial identities, to the primes lying in a small number of residue classes
modulo $840$ that are quadratic-residue classes and hence not amenable to identity-based
coverage: the six classes $1,121,169,289,361,529 \pmod{840}$, the smallest representative being
$p=1009$. The conjecture is exactly equivalent (Bloom's covering-equivalence theorem, with
ancestry in Nakayama, Rosati and Mordell) to an infinite covering-system statement: every
$n$ lies in a residue class $-a/c \pmod{4acd-1}$ for some positive integers $a,c,d$, or in one
of finitely many auxiliary Type-II classes. Equivalently, writing $q=4acd-1$, solvability at $n$
is the existence of a divisor $q\mid cn+a$ with $q\equiv -1\pmod{4ac d}$. The entire difficulty
of the conjecture is concentrated in this covering statement for a density-zero exceptional set:
Vaughan proved in 1970 that the number of $n\le x$ for which $4/n$ is \emph{not} a sum of three
unit fractions is $O\!\left(x\exp(-c(\log x)^{2/3})\right)$, and no exponent has been improved
upon since.

\subsection{What this archive is, and what it is not}
\label{subsec:whatitis}

The archive under report, \file{erdos-straus-main}, is the public git repository of an
AI-assisted mathematical research programme conducted under the auspices of Riley Betts Ltd,
released under the MIT licence (with a single, separately attributed Apache~2.0 exception,
\Cref{subsec:provenance}). Its own front-door document,
\texttt{erdos-\allowbreak straus-\allowbreak programme.md}, states plainly:
\begin{quote}
\small\itshape
``This file is the programme brief for the git repository. It is not a theorem. \textbf{No proof
of the Erd\H{o}s--Straus conjecture is claimed or achieved.} The archive is kept so that the
mathematical community can use the record.''
\end{quote}
This report inherits that discipline. Nothing below should be read as a claim that the
Erd\H{o}s--Straus conjecture has been proved, reduced to a solved sub-problem, or even reduced to
a single named open conjecture that the community regards as tractable. What the archive
contains is:
\begin{enumerate}[leftmargin=1.6em]
  \item a genuine, if partial, \textbf{formal (Lean~4) development}, described in
        \Cref{sec:lean}, that (i) is axiom-audited, (ii) proves the elementary reduction of ES to
        an explicit covering system and a single named analytic interface
        (\lean{AnalyticSurvivorBound}), (iii) formalises a non-trivial slice of the
        Bright--Loughran Brauer--Manin theory for the associated affine cubic surfaces, and
        (iv) verifies the conjecture unconditionally for every $n<100{,}000$ inside the kernel;
  \item a sequence of \textbf{analytic partial theorems and precisely stated negative results}
        on a covering-density reformulation of the conjecture (\Cref{sec:route1}), including a
        later-withdrawn covering-box claim that would have improved on Vaughan's 1970
        exceptional-set bound (Vaughan remains the published comparison; see the dated notice
        above and \Cref{subsec:epower}), together with
        several careful ``look and reject'' investigations of recent analytic number theory
        (Zheng's simultaneous-progressions theorem, Heath-Brown's $\delta$-method,
        Fouvry--Kowalski--Michel--Sawin monodromy independence, and a Maynard-type near-BV slab)
        against the specific obstruction the programme identifies (the \emph{joint
        well-factorable CRT residue} at the Bombieri--Vinogradov wall), culminating in a
        capstone document that diagnoses the common failure mode of all four;
  \item a \textbf{geometric investigation} (\Cref{sec:route2}) of the Bright--Loughran
        Brauer--Manin obstruction theory for the Erd\H{o}s--Straus surfaces, formalising several
        of their theorems, testing (and refuting) three candidate geometric proof mechanisms —
        Harpaz's descent-fibration theorem, a real-compatible twist construction, and a
        triangle-boundary Campana/semi-integral existence principle — and arriving at the
        structural finding that every geometric mechanism considered reduces to the same
        analytic covering-density core as Route~1;
  \item a \textbf{triage of a sibling research effort and an external claimed proof}
        (\Cref{sec:prior}): a prior, separate covering-density assault is examined and mostly
        left out as a recorded negative, keeping only a ported Mathlib-style library theorem and
        a small Type-II identity dictionary; and a 2026 preprint claiming an elementary proof of
        ES is checked directly against the classical Mordell obstruction and found not to
        establish the covering claim it needs;
  \item a substantial body of \textbf{numerical experimentation} in Python
        (\Cref{sec:computation}), instrumenting every quantitative claim above with reproducible
        measurements at scales up to $x=10^9$;
  \item and a disciplined \textbf{research log} (the twenty-three sections \S4a--\S4w of
        \texttt{erdos-\allowbreak straus-\allowbreak novel-\allowbreak structures-\allowbreak plan.md}, together with a chain of dated
        ``postscript'' entries in the programme brief) that records not only positive results
        but every investigated route that was tried and found wanting, with the specific
        obstruction named in each case.
\end{enumerate}
A recurring stylistic feature of the archive, preserved throughout this report, is the sharp
distinction between three categories of statement: results that are \emph{proved} inside the
Lean kernel (flagged \LEAN{} below); results that are \emph{claimed} on paper with a written
mathematical argument but not machine-checked (flagged \CLAIMED{}); and quantities that are
\emph{measured} computationally at finite, stated scales, which is evidence but not proof
(flagged \MEASURED{}). Distinct again are the several \emph{negative} findings, where a
specific published theorem or technique is checked against the programme's target and found not
to apply, for a stated structural reason (flagged \NEGATIVE{}). This taxonomy is not a
stylistic ornament: the archive's central methodological claim is that most of the historical
difficulty of the Erd\H{o}s--Straus conjecture is explained by a single \emph{k-budget
invariant} (\Cref{subsec:kbudget}) that every one of these negative results instantiates in a
different technical guise.

\subsection{Authorship and provenance}
\label{subsec:provenance}

The public archive is explicit about how the work was produced, and this report follows it in
recording that explanation rather than paraphrasing it away. The programme brief states:
\begin{quote}
\small\itshape
``Martyn Riley, of Riley Betts Ltd, facilitates the automation and the high-level strategy
(novel routes, recorded negative results, documentation procedure for the agents). He makes no
claim to expertise in number theory, combinatorics, or geometry. \dots{} The mathematical and
formalization work is a collaboration of Anthropic agents, Grok~4.6, and Cursor agents, with
Lean~4, Mathlib, and numerical test resources. It is not a record of human coauthorship of the
mathematics.''
\end{quote}
This report adopts the same framing throughout: where the text below refers to ``the programme''
producing a result, this denotes the automated research process described above, not a claim of
individual human mathematical authorship. The repository is released under the MIT licence,
copyright Riley Betts Ltd, with one stated exception: \lean{Stormer.lean}
(\Cref{subsec:stormer}) is adapted from a public Mathlib pull request and remains Apache
License~2.0, copyright Alexander Benjamin Worth Burns and Ravi Bajaj. A GitHub Actions workflow
(\file{.github/workflows/ci.yml}) rebuilds the pinned Mathlib-dependent layer on every push, and
a \file{CONTRIBUTING.md} records the archive's pinning policy (Lean toolchain, Mathlib revision,
lockfile, Python numerics) and reiterates the same scope restrictions as the mathematical
documents themselves: no covering densification, no treatment of the external claims of
\Cref{sec:prior} as resolved, no treatment of the library-only files as progress toward the
load-bearing interface.

\subsection{Objectives of the programme}

Stated explicitly across the roadmap document \texttt{erdos-\allowbreak straus-\allowbreak road-\allowbreak to-\allowbreak lean-\allowbreak qed.md} and the
acceptance specification \texttt{erdos-\allowbreak straus-\allowbreak conjecture-\allowbreak spec.md}, the programme's objectives are:
\begin{enumerate}[leftmargin=1.6em]
  \item to convert the folklore statement of ES into a machine-checked reduction whose only
        open content is a single, explicitly named analytic (or geometric) hypothesis;
  \item to identify, as precisely as possible, \emph{why} the conjecture has resisted proof for
        over 75 years — not merely to assert difficulty, but to locate a specific
        information-theoretic ceiling (the k-budget invariant) that every interval-intrinsic
        method is shown to respect, and to enumerate the input classes that are \emph{not}
        capped by that ceiling;
  \item to produce, along the way, genuine unconditional mathematics that stands on its own
        merit regardless of whether the summit is reached — in particular the \(d=1\) floor
        E\_lane, below Vaughan (1970). The covering-box power saving that would have improved
        on that record is \WITHDRAWN{} (\Cref{subsec:epower}); Vaughan remains the published
        comparison;
  \item to formulate an \textbf{acceptance specification} (eight requirements, R1--R8,
        \Cref{sec:spec}) that any candidate ``load-bearing'' conjecture $C$ must satisfy to
        count as a genuine reduction of ES, and to test a wide slate of candidate conjectures
        against it (\Cref{sec:candidates}), several of which are refuted or frozen along the
        way;
  \item and to record, as open mathematical questions addressed to the relevant published
        literatures rather than as correspondence with named individuals, a bounded set of
        questions in arithmetic geometry (integral points on log K3 surfaces,
        \Cref{subsec:loughran-questions}) and in analytic number theory (the half-dimensional/
        vector-sieve literature, \Cref{subsec:sievedesk}) that the programme's own tools cannot
        resolve.
\end{enumerate}

\subsection{Structure of this report}

\Cref{sec:covering} recalls the covering-system reformulation of ES that underlies the entire
programme and the classical reductions it is built on. \Cref{sec:lean} gives a detailed,
file-by-file account of the Lean~4 formalisation. \Cref{sec:route1} discusses the analytic
covering-density route (``Route~1''), including the k-budget invariant, the sequence of
theorems T(1)/T(3)/T(A)/T(A)$^+$/E\_lane/E\_power, the negative results obtained while
attempting to complete the T(3) lower bound, and their capstone diagnosis. \Cref{sec:route2}
discusses the geometric Bright--Loughran route (``Route~2''), including the formalised
Brauer--Manin theory, the Harpaz/Cao--Xu/twist/Campana investigations, the Schinzel separator
programme, and the synthesis that the two routes are one route. \Cref{sec:prior} describes the
triage of the prior Track-1 archive and the examination of an external claimed proof.
\Cref{sec:candidates} catalogues the eleven-plus named candidate conjectures the programme
considered as the load-bearing statement $C$. \Cref{sec:spec} states the acceptance
specification R1--R8 in full and gives the current scorecard. \Cref{sec:computation} describes
the twenty Python scripts and their numerical findings in detail. \Cref{sec:discussion} gives an
overall assessment. \Cref{app:files} is a complete catalogue of every file in the archive. The
full bibliography, drawn from \file{README.md} and cross-referenced throughout, is given at the
end.


\section{The covering-system reformulation}
\label{sec:covering}

\subsection{Witnesses and the covering criterion}

The programme's entire formal and analytic apparatus is built on a single elementary
re-derivation, carried out independently in Lean (\Cref{sec:lean}) and used throughout the
analytic work. Suppose $n$ admits parameters $a,c,d,m\in\Z_{>0}$ with
\[
cn+a+m = 4acdm. \tag{$\ast$}
\]
Call such a quadruple a \emph{witness} for $n$, written $\mathrm{Witness}(n,a,c,d,m)$. A short
computation (the identity is proved as \lean{witness\_sound} in \file{ErdosStraus.lean};
\Cref{subsec:erdosstraus-lean}) shows that a witness yields an explicit solution
\[
4/n \;=\; \frac{1}{adm} + \frac{1}{n\!\cdot\! acd} + \frac{1}{n\!\cdot\! cdm}.
\]
Writing $q=4acd-1$, equation $(\ast)$ is equivalent to $q\mid cn+a$ with cofactor $m$; when
$c=1$ this is the classical Type-I congruence $n\equiv -a\pmod{q}$. The \emph{covering system}
at level $A$ consists of every cell $(a,c,d)$ with $1\le a,c\le A$, $1\le d\le 5$; $n$ is
\emph{covered} at level $A$ if some cell in the box divides $cn+a$, and a \emph{survivor}
otherwise. Because a hit always yields a witness (this is exactly
\lean{ES.Covering.covering\_sound}), the entire conjecture reduces to a statement about
survivors:
\[
\text{ES holds for } n \quad\Longleftarrow\quad n \text{ is covered at some level } A.
\]
The classical polynomial identities of Mordell, Aigner, Schinzel and others are simply
particular, fixed, small-parameter witnesses (or one-parameter families of them); they dispose
of every residue class modulo $840$ except the six \emph{hard classes} listed above. The
programme's Lean development formalises this reduction completely (\Cref{subsec:mordell}),
so that the open content of ES is exactly:
\begin{quote}
\itshape every sufficiently large prime in a hard class is covered by the level-$A(x)$ box, for
a suitable growing schedule $A(x)$, and the finitely many small exceptions can be checked by
computation.
\end{quote}
This is the \lean{HardLandingHypothesis} of the Lean development, and the single open interface
the programme names is \lean{AnalyticSurvivorBound} (\Cref{subsec:coveringlayer}).

\subsection{The divisor-form duality}

An equivalent and, for analytic purposes, more convenient coordinate system is the
\emph{divisor form}: $n$ is covered by the cell $(a,d)$ exactly when
\[
q \mid n + 4a^2 d, \qquad q \equiv -1 \pmod{4ad},
\]
for some divisor $q\ge 3$ of $n+4a^2d$. The programme identifies this (\lean{divisor\_form\_sound}
in \file{ErdosStraus.lean}) as a re-coordinatisation of the same covering box, with $c$ playing
the role of the free cofactor. In this form the covering condition is manifestly a statement
about the prime (or divisor) factorisation of a shifted value $n+4a^2d$, which is the natural
language for sieve theory: it is exactly the setting of Iwaniec's half-dimensional sieve and its
relatives (\Cref{sec:route1}). The quantity
\[
\mathrm{ClassRough}(p,a) \;:\Longleftrightarrow\; p+4a^2 \text{ has \emph{no} divisor } q\ge 3
\text{ with } q\equiv -1 \pmod{4a}
\]
(composite divisors included, not only aligned primes) is the Lean predicate that names one
``slice'' of the $d=1$ covering box, and its joint density over a growing family of slices
$a=1,\dots,A$ is the object $S(A,x)$ studied throughout \Cref{sec:route1}.

\subsection{Classical background and the size of the hard classes}

The six hard classes modulo $840$ are exactly the classes that are quadratic residues modulo
each of $8$, $3$, $5$ and $7$ simultaneously (equivalently, quadratic residues modulo $840$),
and it is a theorem of Mordell (via quadratic reciprocity) that no fixed polynomial witness
family can cover a quadratic-residue class: this is why the classical reductions stop exactly
at these six classes, and why the exceptional set of ES --- if non-empty --- is confined to
them. The programme's Lean development reproves this full reduction from scratch
(\Cref{subsec:mordell}); its analytic and geometric investigations are conducted entirely on
this hard-class population.


\section{The Lean~4 formalisation}
\label{sec:lean}

\subsection{Architecture: two orthogonal layer systems}
\label{subsec:architecture}

The archive contains thirteen Lean source files together with a \file{lakefile.toml},
\file{lean-toolchain} and \file{lake-manifest.json} defining the build, plus a
\file{.github/workflows/ci.yml} that rebuilds the pinned layer on every push. Eleven files sit
at the repository root and are the active Lake library; two further files sit in an
\file{archive/} subdirectory and are explicitly excluded from the build
(\Cref{subsec:legacy}). Two \emph{different} stratifications named ``Layer A/B(/C)'' occur in
the documentation and must not be conflated (the archive itself flags this risk explicitly in
\texttt{erdos-\allowbreak straus-\allowbreak novel-\allowbreak structures-\allowbreak plan.md},
\S4s):

\begin{description}[leftmargin=1.4em]
  \item[File-map layers (\file{README.md}).] \emph{Layer A} is the bare-Lean core —
        \file{ErdosStraus.lean}, \file{BrightLoughran.lean}, \file{ConicFiber.lean},
        \file{TwistDescent.lean}, \file{SchinzelSep.lean}, \file{SchinzelDecide.lean},
        \file{NoVieta.lean}, \file{EPower.lean} — compiling under Lean~$\ge 4.33.0$ with \emph{no imports at all},
        checkable with the bare \texttt{lean} binary and no \texttt{lake}/Mathlib dependency.
        \emph{Layer B} is \file{ErdosStrausQR.lean}, \file{ErdosStrausBLRoute.lean}, and the two
        library-only additions of this revision, \lean{Leochlon.lean} and \lean{Stormer.lean}
        (\Cref{subsec:leochlon,subsec:stormer}), all of which \texttt{import} Mathlib. The first
        two discharge the single deep mathematical input that Layer A leaves as an explicit
        interface (Hilbert/quadratic reciprocity); the latter two are unrelated library material
        kept for documentation reasons (\Cref{sec:prior}) and carry no weight on the proof
        architecture.
  \item[Roadmap layers (\texttt{erdos-\allowbreak straus-\allowbreak road-\allowbreak to-\allowbreak lean-\allowbreak qed.md}, \S14).] A different,
        orthogonal three-way split of the \emph{proof architecture}: \emph{Layer A} (formal
        algebra and the reduction to \lean{HardLandingHypothesis}, fully compiled); \emph{Layer
        B} (the still-open deep analytic or geometric theorem); \emph{Layer C} (the finite
        certificate closing the range below the analytic threshold $X_0$, whose lane is chosen
        among kernel \texttt{decide}, \texttt{native\_decide}, or a future zero-knowledge-attested
        computation).
\end{description}
Both stratifications are in force simultaneously and are used side by side in the archive; this
report follows the source in keeping them distinct, and uses ``file-map Layer A/B'' and
``roadmap Layer A/B/C'' where ambiguity could arise.

\subsection{Build configuration and licensing}

The Layer-B files are built with \texttt{lake} against a pinned revision of \texttt{mathlib4}
(commit \texttt{1a0ef26d9624d64a1ad11853d17c30b8c32f2a10}) under
\texttt{leanprover/lean4:v4.34.0-rc1} (\file{lean-toolchain}, explicitly documented as a release
candidate rather than a moving nightly), declared in \file{lakefile.toml} (package
\texttt{ErdosStraus}, eleven library roots). The Layer-A files require no such dependency and
can be checked with a bare \texttt{lean} invocation per file, as documented in \file{README.md}.
The lockfile \file{lake-manifest.json} is committed as part of the archive, and
\file{CONTRIBUTING.md} directs that it be regenerated and committed alongside any toolchain or
Mathlib bump; \file{.github/workflows/ci.yml} runs \texttt{lake exe cache get} and
\texttt{lake build} on every push to guard against silent drift. Numeric scripts are pinned
separately via \file{requirements.txt} (a single \texttt{numpy} version constraint). The
repository is released under the MIT licence (\file{LICENSE}), copyright Riley Betts Ltd, with
one exception: \lean{Stormer.lean} (\Cref{subsec:stormer}) carries its own Apache License~2.0
notice (\file{LICENSE-APACHE}), copyright Alexander Benjamin Worth Burns and Ravi Bajaj, as
adapted from a public Mathlib pull request.

\subsection{Axiom discipline}
\label{subsec:axioms}

Every theorem in the file-map Layer-A core is accompanied, in \file{README.md}, by an explicit
axiom audit obtained from Lean's \texttt{\#print axioms}. Three axiom profiles occur:
\begin{itemize}[leftmargin=1.4em]
  \item \textbf{Axiom-free} (\texttt{decide}): theorems whose truth is certified by the Lean
        kernel's own evaluator with no additional axiom, e.g.\ \lean{primes\_below\_100\_covered\_kernel},
        \lean{lemma38\_check\_true}, \lean{TwistDescent.descent}, \lean{SchinzelSep.signed\_9\_5},
        \lean{SchinzelDecide.sep\_9\_5}, \lean{SchinzelDecide.sep\_13\_7}, \lean{SchinzelDecide.pos\_4\_5}.
  \item \textbf{Standard} (\texttt{propext}, \texttt{Quot.sound}, and sometimes
        \texttt{Classical.choice}): the great majority of theorems, including the main identity
        \lean{witness\_sound}, the conditional-QED assembly \lean{conditional\_qed\_hard}, the
        conic-fibre correspondence \lean{ConicFiber.fiber\_to\_divisor}/\lean{divisor\_to\_fiber},
        and \lean{NoVieta.es\_unique\_x}.
  \item \textbf{\texttt{native\_decide}} (adds \texttt{Lean.ofReduceBool}): used only for the
        large finite verifications that would be impractically slow under kernel reduction —
        \lean{primes\_below\_10000\_covered}, \lean{hard\_primes\_below\_10000\_covered},
        \lean{okay\_below\_100000} (the verified frontier to $n<100{,}000$), and the
        Brauer--Manin instance checks \lean{thm12\_instance\_1009},
        \lean{witness\_fingerprint\_1009} in \file{BrightLoughran.lean}.
\end{itemize}
The file-map Layer-B files, including the two library additions of this revision, necessarily
carry Mathlib's own axiom footprint (in particular \texttt{Classical.choice} and
\texttt{propext} pervasively, and \texttt{Quot.sound} where quotient types are used in Mathlib's
own development) and are not separately re-audited in \file{README.md}; the table in
\Cref{app:axioms} reproduces the Layer-A audit exactly as published. The roadmap explicitly
reserves a third, not-yet-implemented lane for closing the \texttt{native\_decide} gap at
larger scale: replacing its trusted-compilation axiom with a zero-knowledge-attested execution
proof of the same checker (the lean-tee/SP1 pipeline referenced in
\texttt{erdos-\allowbreak straus-\allowbreak novel-\allowbreak structures-\allowbreak plan.md}, Option~7/Track~C), which is outside
the scope of this archive but is named as the intended mechanism for scaling the verified
frontier from $10^5$ toward the classical computational record of $10^{17}$ (Salez,
arXiv:1406.6307).


\subsection{\file{ErdosStraus.lean} — the core reduction and covering system}
\label{subsec:erdosstraus-lean}

This is the largest and most load-bearing file in the archive (3420 lines). It is organised
into five successive strata, each building on the last.

\subsubsection{Witness algebra and the finite verified frontier}

\begin{itemize}[leftmargin=1.4em]
  \item \lean{IsES n x y z} is the cleared-denominator predicate defined in \Cref{sec:covering}.
  \item \lean{Witness n a c d m} is the covering condition $(\ast)$.
  \item \LEAN{} \lean{witness\_sound}: a witness for $n$ yields the explicit ES solution
        $(x,y,z)=(adm,\,n\!\cdot\! acd,\,n\!\cdot\! cdm)$. This is the main algebraic identity
        the whole programme is built on, proved by a single \texttt{omega}-closed calculation
        after clearing denominators.
  \item \LEAN{} \lean{scale}: a solution for $p$ lifts to any multiple of $p$ — the standard
        observation that it suffices to prove ES for primes.
  \item \LEAN{} \lean{exists\_prime\_factor}: every $n\ge 2$ has a prime factor, proved
        self-containedly (no import of \texttt{Nat.exists\_prime\_and\_dvd}), by strong
        induction with an explicitly defined \lean{IsPrime} predicate and a proof of
        \lean{IsPrime.dvd\_or\_dvd}.
  \item \LEAN{} \lean{LandingHypothesis}: every prime has a witness. \LEAN{}
        \lean{conditional\_qed : LandingHypothesis → ErdosStraus} assembles the full conjecture
        from this one open hypothesis, via \lean{witness\_sound} and \lean{scale}.
  \item \MEASURED (\texttt{native\_decide})~\LEAN{} instances: every prime, and separately every
        hard-class prime, below $10{,}000$ has a witness with parameters $\le 25$; every prime
        below $100$ is verified with the axiom-free kernel \texttt{decide}. A worked example,
        \lean{es\_1009 : IsES 1009\ 276\ 3027\ 92828}, is derived from the witness
        $(a,c,d,m)=(3,1,1,92)$.
  \item An additional one-parameter family, the \emph{one-slot} construction
        (\lean{OneSlot}, Elsholtz--Tao Type~I, complementary to the two-slot \lean{Witness}),
        is proved sound (\lean{oneSlot\_sound}, \lean{oneSlot\_es}) and later folded into a
        \emph{hybrid} covering predicate (\lean{HybridCovered}) that is provably weaker to
        satisfy than the two-slot box alone.
\end{itemize}

\subsubsection{The Mordell reduction to the six hard classes}
\label{subsec:mordell}

A second stratum reproves, entirely inside Lean and by explicit one-line \texttt{omega} proofs
from named polynomial witnesses, the classical elementary reduction:
\begin{itemize}[leftmargin=1.4em]
  \item \LEAN{} \lean{es\_of\_not\_one\_mod\_four}: ES holds outright for every $n\ge 2$ with
        $n\not\equiv 1\pmod 4$ — the class $n\equiv 3\pmod 4$ via witness $(1,2,t+1,1)$, and all
        even $n$ by scaling the $n=2$ solution.
  \item Seven further named polynomial witnesses discharge $n\equiv 2\pmod 3$;
        $n\equiv 5\pmod 8$; $n\equiv 3,5,6\pmod 7$; $n\equiv 7,13\pmod{15}$
        (\lean{w\_two\_mod\_three}, \lean{w\_five\_mod\_eight}, \lean{w\_three\_mod\_seven}, etc.),
        each a single-line \texttt{omega} discharge of the witness equation.
  \item These are combined, via a \lean{HardClass} predicate and a Presburger CRT lemma
        \lean{mod840\_hard} (again closed by \texttt{omega} alone), into
        \LEAN{} \lean{es\_prime\_not\_hard}: \emph{every prime outside the six hard classes
        satisfies ES unconditionally}.
  \item \LEAN{} \lean{HardLandingHypothesis} and \lean{conditional\_qed\_hard} then restate the
        conditional QED so that the only remaining open hypothesis concerns primes
        $\equiv 1,121,169,289,361,529\pmod{840}$ — exactly matching the classical state of the
        art, now machine-checked end to end.
\end{itemize}

\subsubsection{The verified frontier below $100{,}000$}

\LEAN{} \lean{es\_below\_100000}, of type
$2\le n \to n < 100000 \to \exists\, x\,y\,z,\ \mathrm{IsES}\ n\ x\ y\ z$,
is proved by strong induction combined with a self-certifying range check \lean{okay}: at each
$n$, either a proper factor is found (and the solution is obtained recursively by scaling), or a
bounded covering witness is found directly. Because \lean{okay} is self-certifying, no separate
primality-completeness proof is required. This theorem uses the \texttt{native\_decide} axiom
lane described in \Cref{subsec:axioms} and is identified in the roadmap as the seed of the
zero-knowledge-attested Track~C: replacing \lean{ofReduceBool} by a verified execution proof of
the same checker, and scaling the bound from $10^5$ towards $10^{17}$, is described (but not
implemented in this archive) as the intended lean-tee/SP1 pipeline.

\subsubsection{The covering-system layer and the analytic interfaces}
\label{subsec:coveringlayer}

The fourth stratum (\lean{namespace ES.Covering}, roadmap Steps~1--2) formalises the covering
system of \Cref{sec:covering} as executable Booleans and proves everything about it that does
\emph{not} require analytic input:
\begin{itemize}[leftmargin=1.4em]
  \item \lean{coversB a c d n}, \lean{coveredB A n} (a decidable Boolean scan of the level-$A$
        box), \lean{Covered A n}, \lean{Survivor A n} (with decidability instances).
  \item \LEAN{} \lean{covering\_sound}: \lean{Covered A n} implies the existence of a witness —
        the elementary half of the reduction, proved unconditionally for every level $A$.
  \item \LEAN{} \lean{survivor\_typeI\_aps}: a survivor avoids every Type-I arithmetic
        progression in the box — the ``survivor bundle'' used throughout the analytic
        discussion.
  \item \lean{survivorCount A lo hi} and \LEAN{} \lean{no\_survivors\_of\_card\_eq\_zero}: the
        \emph{cardinality-to-QED bridge} — an explicit count of zero survivors in an interval
        kills every survivor there, by simple integrality.
  \item \textbf{The single load-bearing unproven target}, stated but deliberately left as a
        bare \lean{Prop} (roadmap Step~2, ``statement-first discipline''):
        \[
        \texttt{AnalyticSurvivorBound (Alevel : Nat → Nat) (X0 : Nat) : Prop :=}
        \]
        \[
        \forall\, p,\ X_0\le p \to \mathrm{IsPrime}\ p \to \mathrm{HardClass}\ p \to
        \mathrm{Covered}\ (\mathrm{Alevel}\ p)\ p.
        \]
        \LEAN{} \lean{hard\_landing\_of\_interface} and \lean{erdos\_straus\_of\_interface}
        assemble this interface, together with a finite closure hypothesis below $X_0$, into
        the full conjecture. \emph{This definition is the precise formal target that
        Route~1 (\Cref{sec:route1}) attempts, and fails within the scope of this archive, to
        discharge}; \Cref{sec:route1} explains exactly how far short of it the analytic
        results fall, and why.
  \item A parallel \emph{divisor-form} interface, \lean{DivisorLandingBound}, and the kernel
        equivalence \lean{divisor\_form\_sound} (``the divisor criterion IS the covering box's
        free-$c$ slot'', \texttt{erdos-\allowbreak straus-\allowbreak novel-\allowbreak structures-\allowbreak plan.md} \S4q), together with a
        family of bridge theorems (\lean{dfi\_dictionary\_forward}, \lean{dfi\_dictionary},
        \lean{d1\_landing}) connecting the two coordinate systems.
\end{itemize}

\subsubsection{ClassRough, real characters, and the certificate architecture}
\label{subsec:classrough-lean}

The final and most intricate stratum of \file{ErdosStraus.lean} develops, entirely inside Lean,
the \emph{character-certificate} theory that underlies T(A) and T(3) (\Cref{sec:route1}):
\begin{itemize}[leftmargin=1.4em]
  \item \lean{ClassRough p a}: $p+4a^2$ has \emph{no} divisor $q\ge 3$ (composite divisors
        included) with $q\equiv-1\pmod{4a}$.
  \item A hand-rolled theory of \emph{odd real characters} modulo $4,8,12$ is built from
        scratch as completely multiplicative, periodic, $\pm1$-valued functions: \lean{chi4}
        (the character mod $4$, $\chi_4(3)=-1$), \lean{chi8}, \lean{chi4chi8}, \lean{chi3mod12},
        \lean{chi4chi3}, each proved multiplicative (\lean{chi4\_mul}, \lean{chi8\_mul},
        \lean{chi3mod12\_mul}, \dots) and periodic, and packaged as instances of a generic
        \lean{OddRealChar} structure.
  \item \LEAN{} \lean{classRough\_of\_certifies} (the ``certificate direction''): if a real odd
        character $\chi$ modulo $4a$ certifies $p+4a^2$ — every prime factor of $p+4a^2$
        coprime to $4a$ lies in $\ker\chi$ — then $\mathrm{ClassRough}(p,a)$ holds. This is
        proved for general odd real characters via complete multiplicativity
        (\lean{chi\_eq\_one\_of\_dvd}).
  \item The \emph{converse} is proved as an exact biconditional at $a=1,2,3$:
        \LEAN{} \lean{classRough\_iff\_chi4} ($a=1$), \lean{classRough\_a2\_iff\_certificates}
        ($a=2$, two characters $\chi_4,\chi_4\chi_8$ mod $8$), and
        \lean{classRough\_a3\_iff\_certificates} ($a=3$, two characters $\chi_4,\chi_4\chi_3$ mod
        $12$), combined into the single \textbf{T(3) kernel theorem}
        \[
        \textbf{\lean{classRough\_123\_iff\_certificates}}
        \]
        stating that $\mathrm{ClassRough}$ jointly at $a=1,2,3$ is \emph{exactly} the $1\times
        2\times 2$ union of the four resulting certificate assignments. This single theorem is
        the formal anchor of the entire T(3) analytic argument in \Cref{subsec:t3}.
  \item A companion elementary stratum, \lean{ThreeModFourFree} and the theorems
        \lean{exists\_three\_mod\_four\_prime\_factor}, \lean{d1\_a1\_of\_three\_mod\_four\_factor},
        \lean{not\_d1\_a1\_of\_free}, identifies the $a=1$, $d=1$ escapee condition with
        ``$p+4$ has no prime factor $\equiv 3\pmod 4$'' — the classical sum-of-two-squares
        condition, and the base case that calibrates the whole certificate hierarchy.
\end{itemize}

\subsubsection{The Henriot/Nair--Tenenbaum interface: an elementary $d=1$ sieve, in Lean}
\label{subsec:henriot-lean}

The final $\sim 800$ lines of the file (from \lean{coprime\_four\_of\_odd} to
\lean{D1UniformNtBound} at the file's end) carry out, entirely by elementary \texttt{omega}- and
divisibility-level reasoning, a self-contained combinatorial sieve analysis of the $d=1$,
$a=1$ lane — what the documentation calls ``door (b)'' of the k-budget investigation
(\Cref{subsec:kbudget}). This is the only part of the archive's analytic reasoning that is
carried out \emph{inside} the Lean kernel rather than on paper. Its architecture:
\begin{itemize}[leftmargin=1.4em]
  \item An \emph{inner/outer slot} dichotomy (\lean{D1InnerSlot}, \lean{D1OuterSlot},
        \lean{inner\_or\_outer}) partitions the aligned divisors of $p+4$ by size relative to a
        threshold $2D$.
  \item A sequence of purely combinatorial ``hit-spacing'' lemmas
        (\lean{no\_three\_hits\_in\_interval\_len\_q}, \lean{at\_most\_two\_hits\_in\_box},
        \lean{odd\_hits\_impossible\_of\_kq\_gt\_D}, \lean{unique\_large\_aligned\_landing})
        bound how many aligned divisors of $p+4$ can lie in a fixed-length window, by pure
        counting — no analytic input.
  \item A \lean{SmoothTo} predicate and \lean{outer\_landing\_has\_factor\_gt\_two\_D} /
        \lean{outer\_smooth\_has\_hub\_prime} identify the two mutually exclusive escape
        mechanisms for an ``outer'' landing: either a genuinely large prime factor, or
        smoothness through a bounded ``hub'' prime.
  \item This is packaged into \lean{D1NtExceptional}, \lean{D1Remainder}, and finally a
        \emph{named, unproved} target
        \[
        \texttt{D1UniformNtBound (Dlevel : Nat → Nat) (X0 : Nat) (M : Nat → Nat) : Prop}
        \]
        asserting that whenever a Nair--Tenenbaum/Henriot-type mean-value bound $M(p)$ on the
        number of class-rough slots stays strictly below the interval length $D_{\mathrm{level}}(p)$,
        no $d=1$ exceptional escape occurs. \LEAN{} \lean{D1Remainder\_of\_uniform\_nt} and
        \lean{erdos\_straus\_of\_d1\_uniform\_nt} show, unconditionally, that discharging this
        one named hypothesis at the intended schedule $D\sim\exp(c\sqrt{\log x})$ (with $M$ the
        Henriot implicit constant) would complete ES on the $d=1$ lane. The hypothesis itself is
        explicitly \emph{not} discharged in this archive; it is the Lean-level shadow of ``door
        (b)'' in the k-budget investigation (\Cref{subsec:kbudget}), and its status is
        identical to \lean{AnalyticSurvivorBound}: named, wired, open.
\end{itemize}
This stratum is, in effect, a hand-formalised skeleton of a mean-value sieve argument, built to
the point where only one external analytic input (a $k$-uniform Nair--Tenenbaum/Henriot bound)
remains — precisely the ``door (b)'' investigation recorded on paper in
\texttt{erdos-\allowbreak straus-\allowbreak novel-\allowbreak structures-\allowbreak plan.md} \S4g--\S4h (\Cref{subsec:kbudget}).


\subsection{\file{BrightLoughran.lean} — the finite arithmetic heart of Theorem~1.2}
\label{subsec:brightloughran-lean}

This 407-line file formalises the elementary, finite-computation heart of Bright and
Loughran's Theorem~1.2 \cite{BrightLoughran2020}, isolating the single deep external input
(Hilbert reciprocity over $\Q$) as an explicit \lean{structure} interface.
\begin{itemize}[leftmargin=1.4em]
  \item Serre's $2$-adic symbol data (\cite{Serre1973}) is defined directly on residues:
        \lean{eps r} ($=1$ iff $r\equiv 3\pmod 4$) and \lean{omg r} ($=1$ iff
        $r\equiv\pm3\pmod 8$).
  \item \textbf{The finite heart of BL Lemma~3.8.} In reduced $2$-adic coordinates
        $(u_1,u_2,u_3)=(2^{s+e}r_1,2^{s+e}r_2,2^e r_3)$ with $r_i$ odd and $s\ge1$ (a shape
        forced by BL's own valuation analysis), the $2$-adic invariant of the Brauer class
        $\alpha$ is $(-1)^{f+g}$ with $f=\varepsilon(-r_1r_3)\varepsilon(-r_2r_3)$ and
        $g=s(\omega(r_1)+\omega(r_2))$. Because only residues mod $8$ and $(2^s\bmod 8,\,s\bmod
        2)$ matter, the general lemma reduces to a \emph{finite} verification over
        $4^3\times4=256$ cases, executed by \lean{lemma38Check} and proved true —
        \emph{axiom-free} — by \LEAN{} \lean{lemma38\_check\_true : lemma38Check = true := by
        decide}.
  \item \lean{InvariantData}: a \lean{structure} packaging the four local invariants of
        $\alpha$ (at $\infty$, at $2$, at the odd prime $p\mid n$, and the product over all
        remaining good places) together with the single deep field
        \[
        \texttt{reciprocity : invInf * inv2 * invP * invGood = 1}
        \]
        — Hilbert reciprocity, left as an interface on bare Lean and discharged on a machine
        with Mathlib in \file{ErdosStrausQR.lean}.
  \item \LEAN{} \lean{bl\_thm12\_prime\_anatomy}: given $\inv_\infty=-1$ (the positive-octant
        invariant, BL Lemma~3.1), $\inv_2=1$ (odd $n$, BL Lemma~3.8), and $\inv_{\mathrm{good}}=1$
        (BL Lemma~3.5), reciprocity alone forces $\inv_p=-1$ — Yamamoto's non-residue
        condition, i.e.\ exactly the covering-witness fingerprint of \Cref{sec:covering}.
        \lean{yamamoto\_condition} restates this as a statement about the Legendre symbol
        directly.
  \item \MEASURED (\texttt{native\_decide})~\LEAN{}: verified instances of Theorem~1.2 and its
        Yamamoto cross-check at the smallest hard-class prime, $p=1009$:
        \lean{thm12\_instance\_1009} and \lean{witness\_fingerprint\_1009}, tying the geometric
        invariant directly to the covering witness \lean{es\_1009} proved in
        \file{ErdosStraus.lean}.
  \item BL Lemma~3.2 (the \emph{unit-ratio lemma}: among three coordinates with distinct
        $p$-adic valuations, some ratio $u_i/u_j$ is a $p$-adic unit) is proved in full,
        \LEAN{} \lean{lemma32}, together with the supporting Hilbert-symbol machinery
        (\lean{hilbertOdd}, \lean{hilbert2Odd}, \lean{lemma34}, \lean{lemma31}/\lean{lemma31\_pos}
        — the real-place invariant is $-1$ on the positive octant).
\end{itemize}

\subsection{\file{ErdosStrausQR.lean} — the Mathlib discharge}
\label{subsec:erdosstrausqr-lean}

A 170-line file importing \texttt{Mathlib.Tactic} and
\texttt{Mathlib.NumberTheory.\allowbreak LegendreSymbol.\allowbreak QuadraticReciprocity}, this is the smallest
Layer-B file and exists solely to connect the bare-Lean interfaces of
\file{ErdosStraus.lean}/\file{BrightLoughran.lean} to genuine Mathlib number theory.
\begin{itemize}[leftmargin=1.4em]
  \item \LEAN{} \lean{isPrime\_iff\_prime}: the hand-rolled \lean{IsPrime} predicate coincides
        with \lean{Nat.Prime}.
  \item \LEAN{} \lean{typeI\_legendre} / \lean{typeI\_nonresidue\_of\_square\_a}: the
        \emph{Yamamoto fingerprint} — a Type-I witness with prime modulus $q=4ad-1$ forces
        $(p/q)=(-a/q)$ via the Legendre symbol; when $a$ is itself a square and
        $q\equiv3\pmod4$, this forces $(p/q)=-1$, recovering the Schinzel/Mordell Type-I
        obstruction directly from Mathlib's Legendre-symbol API.
  \item \lean{quadratic\_reciprocity\_odd} instantiates Mathlib's
        \lean{legendreSym.quadratic\_reciprocity} for the odd-place content of BL~\S4, and
        \lean{thm12\_instance\_1009\_mathlib}, \lean{witness\_1009\_typeI\_symbol},
        \lean{reciprocity\_1009\_11} reproduce the $p=1009$ verified instances using genuine
        Mathlib reciprocity rather than the bounded-search \lean{legendre} of Layer~A.
  \item \lean{lemma32\_of\_prime} specialises \lean{BL.lemma32} to \lean{Nat.Prime}.
\end{itemize}
Hilbert symbols themselves are not present in Mathlib, so \lean{InvariantData.reciprocity}
remains an interface even in this file; what is discharged here is the odd--odd quadratic
reciprocity content it reduces to, which is the part Mathlib does supply.

\subsection{\file{ErdosStrausBLRoute.lean} — the Bright--Loughran route, in full}
\label{subsec:blroute-lean}

At 3675 lines, this is the single largest file in the archive and the principal deliverable of
the geometric investigation (\Cref{sec:route2}). It develops, on top of the Layer-A/B interfaces
above, the full affine $\Z$-model of the Erd\H{o}s--Straus and Schinzel surfaces, the
Bright--Loughran octant invariant, and an extensive discharge of Hilbert reciprocity for the
specific rational arguments the ES problem produces. Its contents, grouped by theme:

\paragraph{Character algebra and the hybrid covering bypass (plan \S4p--\S4q).}
\lean{witness\_legendre}, \lean{typeI\_d\_slice}, \lean{typeI\_modulus\_11\_pairs},
\lean{hard\_prime\_legendre\_three}, and \LEAN{} \lean{survivor\_legendre\_three} formalise the
finding that Type-I covering exclusions are unions of quadratic-residue \emph{cosets}: a
Legendre-symbol mismatch at a small modulus (e.g.\ $q=11$) excludes an entire $d$-slice at once.
\lean{erdos\_straus\_of\_character\_covering} and \lean{erdos\_straus\_of\_hybrid\_covering}
package this as an alternative sufficient route to ES.

\paragraph{The affine $\Z$-model (BL \S2--\S3).}
\lean{IsSchinzelZ m n x y z} is the integer-point predicate for the affine cubic
$m\cdot xyz=n(xy+yz+zx)$; \lean{IsSchinzelZ 4 n} specialises to \lean{IsESZ}. \LEAN{}
\lean{isES\_to\_esZ}/\lean{isES\_of\_esZ\_pos}: natural-number solutions are exactly the
positive-octant integer points, so \LEAN{}
\[
\texttt{erdos\_straus\_iff\_pos\_octant} : \mathrm{ErdosStraus} \iff \forall n\ge 2,\
U_n(\Z)\cap C_+ \ne \varnothing.
\]
\LEAN{} \lean{esZ\_nonempty}: $U_n(\Z)$ is nonempty for \emph{every} $n\ge 2$ — even $n=2t$ via
the positive point $(2t,2t,t)$, odd $n=2k+1$ via the classical mixed-sign point
$\tfrac1k+\tfrac1{k+1}-\tfrac1{nk(k+1)}$, which lies \emph{off} the positive octant
(\lean{esZ\_of\_odd\_not\_pos}). \lean{schinzelZ\_not\_all\_neg}: no all-negative point exists.

\paragraph{The Jahnel--Schindler/BL box (Lemma~3.10).}
\LEAN{} \lean{schinzelZ\_min\_le}/\lean{esZ\_min\_le}/\lean{es\_pos\_min\_le}: every ordered
integer point satisfies $m\,|u_{\min}|\le 3n$, derived from the factor identity
$(ry-s)(rz-s)=s^2$ with $r=mx-n$, $s=nx$ (\lean{schinzelZ\_factor}, \lean{schinzelZ\_coord\_dvd}).
The converse in the positive octant, \LEAN{} \lean{es\_of\_s2\_divisor}, converts a divisor pair
of $s^2$ at fixed $x$ into a positive point; \emph{existence} of such a pair as $n$ varies is
exactly the conjecture. A single explicit slice ($x=(n+3)/4$, $r=3$) is shown to occupy the
octant for every $n\equiv5\pmod8$ (\lean{es\_five\_mod\_eight\_of\_repair}) and for those
$n\equiv1\pmod4$ whose companion $s$ has a prime factor $\equiv2\pmod3$
(\lean{es\_of\_r3\_slice}); it is shown, honestly, to \emph{fail} at $n=73$
(\lean{r3\_slice\_obstructed\_at\_73}), and the file explicitly does not pursue chaining further
moduli, describing that as ``covering densification''.

\paragraph{The octant invariant and TUB-EP.}
\lean{invInfES} is the real Hilbert symbol of $\alpha$, cleared of the positive square $u_3^2$.
\LEAN{} \lean{invInfES\_eq\_neg\_one\_iff}: it equals $-1$ exactly on the same-sign octants (the
all-negative octant being empty, this detects the positive component — BL Lemma~3.1 on the
$\Z$-model). \LEAN{} \lean{invInfES\_of\_esZ\_mixed}/\lean{invInfES\_odd\_mixed}: mixed-sign
integer points have invariant $+1$ (BL Theorem~1.5 at the real place). The central named
hypothesis of the geometric route is then
\[
\textbf{\lean{TubEpHypothesis}} : \forall n\ge 2,\ \exists\, u_1,u_2,u_3>0,\ \mathrm{IsSchinzelZ}\ 4\ n\ u_1\ u_2\ u_3,
\]
with \LEAN{} \lean{erdos\_straus\_of\_tub\_ep} proved and \LEAN{}
\[
\textbf{\lean{tub\_ep\_iff\_erdos\_straus}}
\]
establishing that \lean{TubEpHypothesis} is \emph{logically equivalent} to the full conjecture
— a fact the report returns to repeatedly in \Cref{sec:route2}, since a genuine reduction must
have independent mathematical content beyond a bare restatement (requirement R2,
\Cref{sec:spec}). The Z-model facts any candidate existence principle may legitimately consume
are packaged as a \lean{structure} \lean{TubEpConsumed} (nonemptiness, the box, the octant
invariant, the mixed-sign-off-$C_+$ fact, and the positive-rational point $(n,n,n/2)$ that shows
directly why Cao--Xu toric strong approximation cannot by itself imply ES).

\paragraph{Hilbert reciprocity, extensively discharged.}
The largest single portion of the file (roughly lines 1250--3000) proves Hilbert-symbol
reciprocity, via Jacobi/$\chi_4$/$\chi_8$ symbols, for every combination of sign and
coprimality that the ES/Schinzel arithmetic actually produces: odd primes, odd coprime
positives, negative odd coprimes, powers of $2$, mixed signs, odd ratios, and specifically the
BL~3.8 shape $(-2^s r_1/r_3,\,-2^s r_2/r_3)$ when the odd parts are pairwise coprime, together
with a shared-odd-factor variant that removes the pairwise-coprimality restriction after
cancelling common squares. Representative named theorems: \lean{hilbert\_reciprocity\_odd\_primes},
\lean{hilbert\_reciprocity\_odd\_coprime}, \lean{hilbert\_reciprocity\_neg\_odd\_coprime},
\lean{hilbert\_reciprocity\_two\_pow}, \lean{hilbert\_reciprocity\_neg\_two\_pow},
\lean{hilbert\_reciprocity\_self}, \lean{hilbert\_reciprocity\_neg\_pos\_odd\_coprime},
\lean{hilbert\_reciprocity\_odd\_ratio}, \lean{hilbert\_reciprocity\_es\_ratio},
\lean{hilbert\_reciprocity\_shared}, \lean{hilbert\_reciprocity\_odd\_integers}. These feed
\LEAN{} \lean{thm12\_discharged} (the bookkeeping $\inv_\infty=-1$, $\inv_2=1$,
$\inv_{\mathrm{good}}=1$, plus reciprocity $\Rightarrow \inv_p=-1$) and
\lean{yamamoto\_of\_nat\_solution}.

\paragraph{Lemma~3.8's two-adic valuation analysis.}
A further self-contained block (lines $\approx3030$--$3600$) proves the two-adic valuation
bookkeeping BL Lemma~3.8 actually needs on \emph{natural-number} solutions rather than the
idealised reduced coordinates of \Cref{subsec:brightloughran-lean}: \lean{es\_v2\_balance},
\lean{exists\_unique\_min\_pair}, \lean{es\_odd\_two\_adic\_vals}, culminating in
\LEAN{} \lean{lemma38\_of\_two\_adic\_shape} and \lean{lemma38\_of\_nat\_solution}, which
derives the reduced $(r_1,r_2,r_3,s,e)$ shape BL~\S3.3.4 requires directly from an arbitrary
positive integer solution, closing the gap between \Cref{subsec:brightloughran-lean}'s finite
lemma and genuine ES solutions. The file closes with \lean{RepresentedX2p4Y2},
\lean{odd\_three\_mod\_four\_free\_is\_form}, and \lean{D1ShiftForm}/\lean{d1\_a1\_iff\_not\_free},
tying the $a=1$ certificate condition of \Cref{subsec:classrough-lean} back to the classical
sum-of-two-squares representation.

\paragraph{What this file does not claim.}
The file's own header is explicit: it does not claim unconditional \lean{ErdosStraus} or
\lean{TubEpHypothesis}; Theorem~1.8 is not formalised and is not treated as existence of a
point; Cao--Xu is not treated as filling the octant; Harpaz is not applicable (\Cref{subsec:conicfiber});
completeness of any covering system is not claimed.


\subsection{\file{ConicFiber.lean} — the split-fibre theorem}
\label{subsec:conicfiber}

This 194-line file is small, entirely self-contained, and proves the single geometric fact that
determines the fate of every fibration-based attack considered in \Cref{sec:route2}. Fixing the
third ES coordinate $t$ in $4txy=n(xy+tx+ty)$ and setting $A=4t-n$, the surface equation is
shown, by pure algebraic manipulation over $\N$ (no case split needed beyond positivity), to be
equivalent to the identity
\[
(Ax-nt)(Ay-nt) = (nt)^2,
\]
proved as \LEAN{} \lean{fiber\_identity}, built from the auxiliary lemmas \lean{four\_t\_gt},
\lean{fiber\_relation} and \lean{coord\_exceeds}. Consequently \emph{every} conic fibre of the
Erd\H{o}s--Straus surface is a \emph{split} conic $uv=c$, for every $n\ge2$ and every fibre $t$
— not merely generically. Two further theorems make this constructive:
\begin{itemize}[leftmargin=1.4em]
  \item \LEAN{} \lean{fiber\_to\_divisor}: a positive point on fibre $t$ yields an explicit
        positive divisor pair $(u,v)$ of $(nt)^2$ with $u\equiv v\equiv -nt\pmod{4t-n}$, and the
        coordinates recover linearly, $(4t-n)x=u+nt$.
  \item \LEAN{} \lean{divisor\_to\_fiber}: conversely, any such positive divisor pair in the
        correct residue class produces a positive point on fibre $t$.
\end{itemize}
Together these prove, unconditionally and for \emph{all} $n\ge2$, that along every conic
fibration of the ES surface, the positive-point existence question \emph{is} the divisor
problem of \Cref{sec:covering} — the geometric and analytic routes are, at the level of this one
theorem, the same object viewed in two coordinate systems. The immediate consequence, discussed
at length in \Cref{subsec:harpaz}, is that Harpaz's descent-fibration theorem for conic log~K3
surfaces \cite{Harpaz2019} cannot apply here: its hypotheses require \emph{non-split} conic
fibres carrying genuine Selmer structure, and none exists on this surface. The file's axiom
profile is the minimal \texttt{propext}, \texttt{Quot.sound} pair throughout.

\subsection{\file{TwistDescent.lean} — the real-compatible twist and its Gate~3 failure}
\label{subsec:twistdescent}

A 219-line file implementing the corrected form of candidate conjecture C5 (\Cref{sec:candidates}):
the \emph{real-compatible twist} $V_{n,d}$, obtained by adjoining $w^2=d\cdot u_1u_3$ ($d>0$
squarefree) to the affine model $U_n$.
\begin{itemize}[leftmargin=1.4em]
  \item \lean{OnUn}, \lean{IsTwist n d x y z w}, \lean{PosRealLocus} set up the objects.
  \item \LEAN{} \lean{naive\_cover\_empty}: the \emph{naive} $\alpha$-trivialising cover
        $w^2=-u_1u_3$ has no real points at all over the positive octant (since
        $u_1,u_3>0\Rightarrow u_1u_3>0$) — the kernel form of the refutation of the original,
        naive version of candidate C5 (\Cref{sec:candidates}).
  \item \textbf{Gate~2} (the cheap half of requirement R1, \Cref{sec:spec}): \LEAN{}
        \lean{descent}: an integral point of $V_{n,d}$ in the positive real locus projects to a
        positive-octant point of $U_n$. \LEAN{} \lean{lift}: conversely, every positive-octant
        point of $U_n$ lifts to some twist (with $d=u_1u_3$, not claimed squarefree-minimal).
  \item \textbf{Gate~3} (\Cref{subsec:c5}, the decisive negative finding): fibering $V_{n,d}$ by
        $u_2=t$, an \emph{explicit rational section} over $\Q(t)$ is exhibited and its cleared
        form proved as a Nat identity, \LEAN{} \lean{t\_fiber\_section\_cleared}. A rational
        section means the generic fibre is \emph{split}, so Harpaz's Theorem~1.0.1 — whose
        engine requires non-split fibres — does not apply to this fibration either, for the
        same structural reason as \Cref{subsec:conicfiber}.
  \item \textbf{Autopsy} (20 August 2026): the very same rational curve, reparametrised, is
        shown to also be a section of the projections $\pi=u_1$ and $\pi=w$
        (\LEAN{} \lean{u1\_fiber\_section\_cleared}, \lean{u1\_fiber\_twist}) — changing the
        linear projection does not produce a genuinely non-split conic fibration of $V_{n,d}$.
\end{itemize}
This file is the kernel record of candidate C5's refutation chain (C5 $\to$ C5$_1$ $\to$
C5$_2$, \Cref{subsec:c5}): a mathematically substantial idea, formally checked, and formally
closed within the same investigation.

\subsection{\file{SchinzelSep.lean} and \file{SchinzelDecide.lean} — the Schinzel separator suite}
\label{subsec:schinzel-lean}

Two short files (84 and 268 lines) implementing the S1/S1$_1$ investigation
(\Cref{subsec:schinzel}) into the generalised Schinzel family $m/n=1/x+1/y+1/z$.
\begin{itemize}[leftmargin=1.4em]
  \item \file{SchinzelSep.lean}: \lean{OnSchinzel m n x y z} is the general cleared model.
        \LEAN{} \lean{signed\_9\_5}: the explicit mixed-sign integer point $(-5,1,1)$ satisfies
        $\tfrac95=1+1-\tfrac15$ (axiom-free \texttt{decide}). \LEAN{} \lean{no\_pos\_9\_5}: for
        \emph{ordered} $x\le y\le z$, no positive solution of $9/5$ exists — proved by a direct
        elementary bounding argument (the leading coefficient is forced to $x=1$, then
        $y\in\{1,2\}$, both residues fail), entirely by hand, no search.
  \item \file{SchinzelDecide.lean}: a genuinely \emph{decidable} procedure
        \lean{decideCplus m n}, searching $x\le 3n$, $y\le 2nx$ and recovering $z$ from
        $nxy=zD$. \LEAN{} \lean{decideCplus\_sound} and, more substantially,
        \LEAN{} \lean{decideCplus\_complete} prove that this bounded search is both sound
        \emph{and complete} for positive-octant solvability of \emph{any} fixed $(m,n)$ — a
        genuine finite decision procedure, proved correct by an explicit chain of bounding
        lemmas (\lean{exists\_sorted}, and bounds \texttt{hb1}--\texttt{hb4} inside the
        completeness proof) rather than by brute enumeration alone.
        \LEAN{} \lean{sep\_9\_5}, \lean{sep\_13\_7} (both axiom-free \texttt{decide}): the
        checker returns \texttt{false} on $9/5$ and on the companion separator $13/7$ (from the
        family $(2z-1)/z$ at $z=7$), yielding \LEAN{} \lean{no\_pos\_9\_5} and
        \lean{no\_pos\_13\_7} as corollaries; \lean{pos\_4\_5} is a positive control
        ($4/5=\tfrac12+\tfrac14+\tfrac1{20}$).
\end{itemize}
These two files jointly establish, unconditionally, that \emph{mixed-sign integer occupancy of
a Schinzel surface does not imply positive-octant occupancy} — the necessity check that drives
the refinement of candidate conjecture C11 (\Cref{claim:c11}).

\subsection{\file{NoVieta.lean} — no Markoff-type involution}
\label{subsec:novieta}

A 156-line file proving the single structural fact that closes the ``geometric mechanism
inventory'' of \Cref{subsec:synthesis}: the affine equation $m\,xyz=n(xy+yz+zx)$ is
\emph{linear} in each coordinate separately (unlike the Markoff equation
$x^2+y^2+z^2=3xyz$, which is quadratic in each variable and hence admits the Vieta involution
$x\mapsto 3yz-x$ underlying the Bourgain--Gamburd--Sarnak strong-approximation machinery
\cite{BGS2016}).
\begin{itemize}[leftmargin=1.4em]
  \item \lean{denomX m n y z} $= myz-ny-nz$ is the fibre denominator for the $x$-coordinate;
        \LEAN{} \lean{linear\_in\_x}: $x\cdot\mathrm{denomX} = n\,yz$ — the equation really is
        linear in $x$ once $(y,z)$ is fixed.
  \item \LEAN{} \lean{unique\_x}: given a nonzero fibre denominator, $x$ is uniquely determined
        by $(y,z)$ — \emph{there is no second root}, hence no Vieta-type involution and no
        Markoff-style correspondence orbit.
  \item \LEAN{} \lean{es\_unique\_x} specialises this to the ES equation ($m=4$) on positive
        naturals.
\end{itemize}
This is a small file but a structurally important one: it is the formal reason, verified rather
than merely asserted, that the entire Bourgain--Gamburd--Sarnak/Markoff-orbifold literature
(\Cref{subsec:synthesis}) cannot be imported wholesale to the Erd\H{o}s--Straus setting.

\subsection{\lean{Leochlon.lean} — a Type-II identity dictionary}
\label{subsec:leochlon}

A 95-line library file, new in this revision, importing \file{ErdosStraus.lean} together with
Mathlib's \lean{Linarith} and \lean{Ring} tactics. It formalises a single Egyptian-fraction
identity reported for an external, informally documented construction (`leochlon/erdstrau'; the
archive notes that the originating URL was not reachable when first recorded, and treats the
construction purely as a mathematical statement to be checked on its own terms):
\[
(4b-1)(4c-1) = 4p\delta+1, \qquad \delta\mid bc.
\]
\begin{itemize}[leftmargin=1.4em]
  \item \lean{LeochlonWitness p delta b c} (parameters $p,\delta,b,c$) packages the hypotheses; \LEAN{}
        \lean{leochlon\_key\_identity} rewrites the defining equation over $\Z$ into the
        cancellation-free form $4bc = b+c+p\delta$.
  \item \LEAN{} \lean{isES\_of\_leochlon}: any such witness always yields the explicit Type-II
        Erd\H{o}s--Straus triple $(pb,\,pc,\,bc/\delta)$ — i.e.\ $p$ divides two of the three
        denominators, exactly the shape of the archive's own \lean{Witness} predicate
        (\Cref{subsec:erdosstraus-lean}).
  \item \LEAN{} \lean{leochlon\_of\_q\_succ\_div\_four} specialises to $b=(q+1)/4$ when
        $q\equiv3\pmod4$, connecting the construction to the archive's own small-modulus
        covering cells.
\end{itemize}
The file's own header, and the companion note \texttt{erdos-\allowbreak straus-\allowbreak leochlon.md}, are explicit
that this is a \emph{dictionary} into \lean{ES.IsES} and nothing more: the identity is not
automatically a covering-landing witness at any particular $(a,c,d,m)$, and it is not treated as
a gateway or finite-catalogue covering construction. \Cref{sec:prior} discusses the provenance
and scope of this file in full.

\subsection{\lean{Stormer.lean} — a ported library theorem, out of scope for the proof}
\label{subsec:stormer}

An 855-line file, new in this revision and licensed separately from the rest of the archive
(\Cref{subsec:provenance}), proving St\o rmer's 1897 theorem \cite{Stormer1897} on consecutive smooth numbers,
adapted from a public Mathlib pull request (\#42040) to build outside Mathlib's own
\texttt{module} root. The file's own header states plainly that it is ``library only'' and
``not on the T(A)/geometry QED path''. Its mathematical content: given a positive solution of
the Pell equation $x^2=1+dy^2$ whose $y$-coordinate has no prime factor outside a fixed finite
set together with the Pell parameter $d$, the solution is forced to be the \emph{fundamental}
one (\LEAN{} \lean{eq\_fundamental\_of\_primeFactors\_y\_subset\_parameter} and
\lean{eq\_of\_primeFactors\_y\_subset\_parameter}, built on an explicit prime-divisibility lemma
for the Lucas-sequence coordinate \lean{pellY}, \LEAN{}
\lean{exists\_prime\_dvd\_pellY\_not\_dvd\_parameter}). This is then packaged into a
three-valued ``exponent code'' on $s$-smooth numbers (\lean{exponentCode},
\lean{consecutiveCode}) and used to bound the number of consecutive pairs of $s$-factored
naturals: \LEAN{} \lean{card\_consecutive\_factoredNumbers\_le\_sub} gives an explicit bound of
$3^{|s|}$ minus a correction, and \LEAN{} \lean{finite\_consecutiveFactoredNumbers} concludes
finiteness of the set of such consecutive pairs for any fixed finite prime set $s$. None of this
material is used anywhere else in the archive; it is included, and separately licensed, purely
as a documented piece of provenance from the sibling investigation discussed in \Cref{sec:prior}.

\subsection{Frozen and legacy files}
\label{subsec:legacy}

Two further Lean files are present but explicitly marked as non-canonical by \file{README.md},
and are kept in an \file{archive/} subdirectory excluded from the \file{lakefile.toml} build:
\texttt{archive/\allowbreak ErdosStraus-\allowbreak core.lean} (2073 lines) is described as ``a frozen snapshot of an earlier
covering file, not a Lake root'', and \texttt{archive/\allowbreak ErdosStraus-\allowbreak mathlib-\allowbreak project.lean} (2480 lines) as
``a duplicate of \file{ErdosStraus.lean}''. Both share the same header and opening definitions
as the canonical \file{ErdosStraus.lean} (\Cref{subsec:erdosstraus-lean}) but represent earlier
development states; a short \texttt{archive/README.md} states explicitly that they are not Lake
roots and should not be imported. This report does not describe them further.

\subsection{Summary of what is, and is not, proved}
\label{subsec:lean-summary}

Collecting the strata above, the unconditional content of the Lean formalisation is:
\begin{itemize}[leftmargin=1.4em]
  \item the full elementary reduction of ES to \lean{HardLandingHypothesis} (six hard classes
        mod 840), with every classical polynomial-witness case machine-checked;
  \item soundness of the covering system at every level (\lean{covering\_sound}) and the
        cardinality-to-QED bridge;
  \item unconditional verification of ES for every $n<100{,}000$;
  \item the finite arithmetic heart of Bright--Loughran's Theorem~1.2 (Lemma~3.8's $256$-case
        table), verified instances at $p=1009$, and an extensive discharge of Hilbert
        reciprocity for the rational arguments the ES problem actually produces;
  \item the $\Z$-model nonemptiness, box bound, and octant-invariant theorems, and the
        equivalence \lean{TubEpHypothesis}~$\iff$~\lean{ErdosStraus};
  \item the split-fibre theorem and its exact divisor correspondence, for every $n\ge2$ and
        every conic fibre;
  \item the Gate~2/Gate~3 twist-descent theorems, including the rational-section proof that
        closes candidate C5;
  \item the Schinzel-separator decision procedure and its verified instances at $9/5$ and
        $13/7$;
  \item the no-Vieta uniqueness theorem;
  \item and, as library material carrying no weight on the reduction, a checked Type-II
        identity dictionary (\lean{Leochlon.lean}) and a ported Mathlib-style theorem on
        consecutive smooth numbers (\lean{Stormer.lean}), both discussed further in
        \Cref{sec:prior}.
\end{itemize}
What remains explicitly, and by design, unproved and named as such are exactly two interfaces:
\lean{AnalyticSurvivorBound} (equivalently \lean{DivisorLandingBound}, equivalently
\lean{D1UniformNtBound} on the restricted $d=1$ lane) on the analytic side, and
\lean{TubEpHypothesis} on the geometric side — with the kernel theorem
\lean{tub\_ep\_iff\_erdos\_straus} establishing that the latter has, at present, no content
beyond a restatement of the conjecture itself. Neither interface is discharged anywhere in this
archive.


\section{Route 1: the analytic covering-density approach}
\label{sec:route1}

\subsection{The research log and the identification of the k-budget invariant}
\label{subsec:kbudget}

Sections \S1--\S4l of \texttt{erdos-\allowbreak straus-\allowbreak novel-\allowbreak structures-\allowbreak plan.md} record a single day (18
August 2026) of rapid, iterated attempts on a \emph{covering density lemma} — an upper bound on
the density of integers avoiding every congruence class in a growing covering box — together
with the honest documentation of where each attempt stalled. The log is unusual for a research
document in that it preserves refuted intermediate claims rather than silently correcting them;
this report follows that discipline and reconstructs the chain in order, since the final
negative result (the k-budget invariant) is only meaningful in light of what it subsumes.

\paragraph{The ladder of reformulations (\S4a--\S4d).}
An initial \emph{Landing Lemma} (roots of the quadratic congruence $\nu^2\equiv -pd\pmod q$
landing in a divisor-constrained position) is identified as equivalent, via the classical
dictionary of quadratic forms and ring class fields, to a Duke--Friedlander--Iwaniec-style
equidistribution statement (\cite{DFI}; this is ``Option~9/Option~5 fusion'' in the plan's own
numbering). A first quantitative version, the \emph{Covering Density Lemma}, conjectures
$u(A)\ll A^{-2-\varepsilon}$ for the density of survivors of the level-$A$ box, calibrated
against a measured constant $\kappa\approx0.139$ (the exponent of $\log^2A$ in the covering
mass). A \emph{partial theorem} is then proved rigorously by restricting to prime cofactor $m$
and using exact CRT independence across distinct primes plus Bombieri--Vinogradov, giving an
unconditional power-law density bound with exponent $\kappa^*\approx0.29$ — the first genuine
theorem of the investigation, and a template for everything that follows.

\paragraph{The Suen conditioning scheme (\S4e).}
Conditioning on all primes $\le T$ turns the survivor-density problem into a product space over
large-prime coordinates, to which \emph{Suen's correlation inequality} \cite{Suen1990,Janson1990}
applies directly:
\[
u \;\le\; \exp\!\bigl(-\mu+\Delta e^{2\delta}\bigr),
\]
with $\mu$ the total conditioned mass, $\Delta$ the mass on dependent pairs, and $\delta$ the
maximum neighbour mass. Measured gates at toy scale ($A=24$--$48$) show that a small-prime
``hub'' (the primes $3,7,11,13$) is the entire source of dependence, and that conditioning it
away reduces $\Delta,\delta$ to $O(\mu^2/T),O(\mu/T)$ — an elementary second-moment bound, later
made fully rigorous as Lemma~SM (\Cref{subsec:epower}). This is identified as, in principle, a
route to the \emph{full} super-polynomial covering-density bound $u(A)\le\exp(-(1-o(1))\kappa\log^2A)$.

\paragraph{The Vaughan barrier and its correction (\S4f--\S4h).}
A first (later corrected) attempt to transfer this density bound into an exceptional-\emph{count}
bound over a growing interval $[x,2x]$ derives a sieve-dimension constraint
$\theta<1/3$ for a covering level $A=\exp((\log x)^\theta)$, and — remarkably — recovers
\emph{exactly} Vaughan's 1970 exponent $\exp(-c(\log x)^{2/3})$ at the boundary $\theta=1/3$,
independently from three unrelated technologies (the fundamental lemma, the large sieve, and
Vaughan's original exponential-sum method). This ``triple convergence'' is initially read as
strong evidence for a genuine barrier. A subsequent, more careful trace of the relevant Nair--
Tenenbaum/Henriot machinery (\S4g--\S4h, ``door (b)''; this is the analysis formalised in
\Cref{subsec:henriot-lean}) discovers that the original barrier arithmetic \emph{conflated total
sifted mass with sieve dimension}. Corrected, the achievable cap is not
$\exp(-c(\log x)^{2/3})$ but a genuine \emph{power saving} $x^{-c\kappa}$, i.e.\ an
exceptional-set bound $O(x^{1-c})$ that would be \textbf{strictly stronger than Vaughan} if
the transfer could be made. This
correction — explicitly flagged as such in the document, with the original \S4f retained
unedited as a record — is the origin of the later-withdrawn E\_power claim
(\Cref{subsec:epower}). Do not read the corrected cap as a compiled improvement on Vaughan.

\paragraph{The k-budget invariant, in final form (\S4k--\S4l).}
The corrected barrier is then re-derived a \emph{third} independent way, via a
circuit-complexity/``fooling'' argument (intervals fooling a depth-$2$ CRT formula, in the
style of Braverman--Bazzi-type constructions), and the three derivations are triangulated into
a single invariant:
\begin{quote}\itshape
\textbf{The k-budget invariant.} The interval-intrinsic independence budget is $k\cdot\log A\le
\log x$, while QED-precision requires effective independence $k\gtrsim\text{mass}>\log x$. Every
method consuming only interval-intrinsic randomness — sieves, the large sieve, Shiu-type
transfers, $k$-wise-independence fooling — is capped by this inequality.
\end{quote}
The invariant is not claimed as a theorem (the plan calls it ``conjectural, proof-shaped''), but
it organises the entire remainder of the programme: it identifies exactly two input classes that
are \emph{not} capped by it — genuinely spectral/bilinear arithmetic input (Kloosterman/DFI-type
cancellation), and algebraic rigidity used non-statistically (the reciprocity/character-coset
structure of \Cref{subsec:t_a}) — and every subsequent ``look'' investigation in
\Cref{subsec:t3-negative} is, explicitly, a test of whether some named modern technique
supplies input from one of these two classes at the specific modulus where the programme's own
argument stalls.

The resulting statement of the residual open problem is named \textbf{H\_ES} (Survivor Level of
Distribution): at level $A=\exp(c\sqrt{\log x})$ with $\kappa c^2>1$, the survivor count in
$[x,2x]$ should satisfy $\#S\ll x\,e^{-(1-o(1))\kappa\log^2A}$. This is, precisely,
\lean{AnalyticSurvivorBound} restated analytically; it is not proved anywhere in the archive,
and the k-budget invariant is the programme's own explanation of why it should be hard.

\subsection{T(A): the joint density of class-roughness at fixed width}
\label{subsec:t_a}

The central working object of the programme's second phase (20--21 August 2026,
\file{erdos-straus-T-A.md}) is a reformulation of the density question in terms of the
\lean{ClassRough} predicate of \Cref{subsec:classrough-lean}. For $p$ ranging over hard primes,
write
\[
\rho_a(x) = \Pr\bigl[\mathrm{ClassRough}(p,a)\bigr], \qquad
S(A,x) = \Pr\bigl[\mathrm{ClassRough}(p,a)\ \text{for all}\ 1\le a\le A\bigr],
\]
and
\[
\widehat C(A,x) \;=\; \frac{S(A,x)}{\prod_{a=1}^A \rho_a(x)}
\]
for the \emph{selection constant}. The conjectural statement T(A) is that, for each fixed $A$,
$S(A,x)\asymp C(A)\prod_a\rho_a(x)$ with $C(A)$ an explicit, trackable sequence, and the
programme's stated exposure is
\[
\limsup_{A\to\infty} \frac{\log C(A)}{\log^2 A} \;<\; \kappa \;=\; 0.139
\]
(the same constant as \Cref{subsec:kbudget}), \emph{together with} the requirement that extra
covering width continue to contribute genuine mass rather than becoming what the archive calls
\emph{dummy covering} — the two ``die conditions'' tracked throughout.

\paragraph{The certificate architecture.} $\mathrm{ClassRough}(p,a)$ is shown (\Cref{subsec:classrough-lean})
to be \emph{certified} by any real odd character $\chi\bmod 4a$ all of whose non-trivial values
on prime factors of $p+4a^2$ avoid the forbidden class — this is the exact, kernel-checked
direction. The converse (every class-rough $p$ has \emph{some} certifying character) is proved
as an exact biconditional only at $a\in\{1,2,3\}$ (and, by a parallel elementary argument, at
$a=6$) — precisely the values of $a$ for which the unit group $(\Z/4a\Z)^\times$ is an
elementary abelian $2$-group. \MEASURED{} At general $a$ the converse is an \emph{$\omega$-gap},
not a near-equivalence: at $x=10^6$, $A\le40$, the certificate covers only $84.6\%$ of
class-rough survivors, with the shortfall growing sharply with the number of coprime prime
factors $\omega(N)$ of the shifted value ($100\%$ at $\omega=1$, falling to $35\%$ at
$\omega=4$); this is the \emph{defect layer}, shown (\file{c4\_defect\_layer.py}) to be exactly
the combinatorial phenomenon that $-1$ can lie in the subgroup generated by the prime residues
of $N$ without being a literal subset product, which is possible precisely when
$\lambda(4a)\nmid 2$, i.e.\ for $a\notin\{1,2,3,6\}$.

\subsection{T(3): a two-sided bracket, upper half claimed}
\label{subsec:t3}

The \textbf{T(3) kernel theorem} (\lean{classRough\_123\_iff\_certificates}) makes $a=1,2,3$
exact: $\mathrm{ClassRough}$ jointly on these three slices is precisely the union of four
certificate assignments built from the characters $\chi_4$ (mod $4$), $\chi_4,\chi_4\chi_8$
(mod $8$), and $\chi_4,\chi_4\chi_3$ (mod $12$). \MEASURED{} At $x=10^6$ this is verified against
$2{,}370$ hard primes with an exact match (656 class-rough on all three slices, 656 in the
certificate union).

\begin{claimstmt}[T(3)$^+$, \CLAIMED] \label{claim:t3plus}
\[
S(3,x) \;\ll\; (\log x)^{-3/2} \qquad (x\to\infty),
\]
with an effective, explicit implied constant.
\end{claimstmt}
The proof (given in full in \file{erdos-straus-T-3.md}) is a genuine, if not machine-checked,
Selberg upper-bound sieve of dimension $3/2$ on the six hard classes, sieved to level
$z=x^{1/4}(\log x)^{-B}$, with Bombieri--Vinogradov supplying the remainder term at the required
level $z^2\le x^{1/2-}$. The dimension $3/2$ arises because each of the three certificate
kernels excludes a density-$1/2$ set of primes $q$, and the three densities sum (up to $O(1)$
collision corrections, verified to be exactly zero for the relevant small moduli) to $3/2$.
\begin{claimstmt}[T(1), \CLAIMED, by citation]
$\mathrm{ClassRough}$ at $a=1$ is exactly the $\chi_4$-certificate, i.e.\ $p=x^2+y^2-4$ for hard
$p\equiv1\pmod4$, and by Iwaniec \cite{Iwaniec1974} together with the congruence-layer extension
of Fuchs--Hsu--Rickards--Schindler--Stange \cite{FHRSS2025}, $S(1,x)\asymp(\log x)^{-1/2}$.
\end{claimstmt}
This one-slice asymptotic calibrates the exponent $\kappa=1/2$ per certificate that the
dimension-count above extrapolates.

\paragraph{The matching lower bound: a three-step plan with a named stall.}
The programme is explicit that the two-sided T(3) bracket is \emph{not} claimed. The plan for
$c_->0$ is: \textbf{(a)} Iwaniec's half-dimensional sieve \cite{Iwaniec1976} gives a genuine
lower bound per certificate slice (dimension $1/2<1$ evades the parity/sifting-limit problem);
\textbf{(b)} the Br\"udern--Fouvry vector-sieve inequality \cite{BrudernFouvry1994,BrudernFouvry1996}
(with Nath--Xie \cite{NathXie2025} as the two-component template) combines lower-bound weights
across the three simultaneous kernels; \textbf{(c)} \emph{completion} — passing from ``no bad
factor below $z$'' to the genuine certificate condition requires controlling factors
$q>x^{1/2}$, which sits beyond Bombieri--Vinogradov and is where the argument can, and the
programme believes does, stall. Two \emph{escapes} are recorded for the eventual paper's
outlook rather than claimed: a weaker $c_-$ (a positive proportion of the expected order, via
``almost-certificates''), or switching to the divisor-sum detector
$r_\chi(n)=\sum_{d\mid n}\chi(d)$ and a dispersion-method (Linnik \cite{Linnik1960}) attack. The
sharp remaining mathematical question, stated precisely and repeatedly across
\file{erdos-straus-T-3.md}, \texttt{erdos-\allowbreak straus-\allowbreak sieve-\allowbreak desk.md} and
\texttt{erdos-\allowbreak straus-\allowbreak novel-\allowbreak structures-\allowbreak plan.md} \S4w, is:
\begin{quote}\itshape
Does a joint well-factorable weight framework exist for a CRT residue
$\alpha(d_1,d_2,d_3)\bmod\operatorname{lcm}(d_1,d_2,d_3)$ that depends on the summed triple, and
if not, can one be built for this specific three-kernel case (conductors $4,8,12$)?
\end{quote}

\subsection{Four negative results at the completion stall}
\label{subsec:t3-negative}

Between 21 August 2026 and the archive's final state, four separate ``look and reject''
investigations were carried out against exactly this stall, each testing a specific recent piece
of analytic machinery. All four are honestly negative, and each pins down precisely why.

\paragraph{(i) Zheng's simultaneous arithmetic progressions \NEGATIVE{} (\file{erdos-straus-T-2.md}, \texttt{zheng-\allowbreak wellfactorable.md}).}
Zongkun Zheng's theorem on primes in two simultaneous arithmetic progressions
\cite{Zheng2025} is tested as a two-modulus ($T(2)$) mean-value input. A first technical
obstacle (\texttt{zheng-\allowbreak wellfactorable.md}) is that the natural weight
$\lambda_d=\chi(d)\mathbf 1_{d\le D}$ is \emph{not} automatically well-factorable in Zheng's
sense (Definition~1.1 is a Dirichlet-convolution statement, not a pointwise factorisation), a
subtlety resolved by isolating the squarefree part. Substantively: after switching, both
divisors $d_1,d_2$ lie in $[0,\sqrt x]$, and Zheng's Theorem~1.1 restricts one leg to
$\theta\le7/36$. \MEASURED{} On a $181\times181$ grid of the switched exponent plane, the
symmetric cell $d_1\approx d_2\approx\sqrt x$ — exactly the cell the completion step needs — lies
in \emph{neither} Bombieri--Vinogradov nor Zheng's Theorem~1.1 or~1.2; the union of ordinary BV
and Zheng covers only $0.518$ of the relevant plane, and an ``uneven'' reroute that leaves one
of the three T(3) kernels in a Type~I/polylog slot and applies Zheng to the remaining pair still
fails on the cell $(0,1/2,1/2)$ (\file{t2\_zheng\_ranges.py}). \textbf{Verdict:} range no-go;
not T(2)$^+$; not T(3) progress.

\paragraph{(ii) Heath-Brown's $\delta$-symbol method \NEGATIVE{} (\texttt{erdos-\allowbreak straus-\allowbreak T-\allowbreak 3-\allowbreak delta.md}).}
Heath-Brown's dissection identity \cite{HeathBrown1996} is tested as a way to detect the moving
CRT congruence jointly rather than sequentially, hoping the Farey/dissection parameter
$Q_\delta$ supplies a genuinely extra degree of freedom. A short lemma
(``Farey--$k$ incompatibility'') shows that at the Bombieri--Vinogradov wall ($M+N=1$,
$Q_D=1/2$), opening a bilinear cell against the dissection modulus forces $Q_\delta>1/2$, while
completing the residual AP-index sum forces $Q_\delta\le1/2$ — the same identity as the Phase~1
stall, merely repackaged (\file{t3\_delta\_ranges.py}). At the specific uneven cell
$(0,1/2,1/2)$ this particular incompatibility is vacuous (the AP index is absent there), but no
positive saving is obtained up to $D\sim x$ or $D\sim x^{3/2}$ either. \textbf{Verdict:} does not
avoid the wall; no range gained; not T(3) progress.

\paragraph{(iii) Joint monodromy independence \NEGATIVE{} (\texttt{erdos-\allowbreak straus-\allowbreak T-\allowbreak 3-\allowbreak monodromy.md}).}
The 2026 Fouvry--Kowalski--Michel--Sawin theorem on bilinear forms with trace functions
\cite{FKMS2026} is tested as a route to square-root cancellation in the product
$\chi_2(d_2)\chi_3(d_3)K(\cdot\,;D)$ for two moving quadratic-character moduli. Four structural
mismatches are identified: the relevant kernel is not a trace function of a single
$\ell$-adic sheaf on a product variety over a \emph{fixed} finite field (the modulus $D$ moves
with the summation point); the individual characters $\chi_2,\chi_3$ are rank-$1$ Kummer sheaves
with abelian monodromy, not the required ``gallant'' (irreducible, quasisimple) type; after
switching, a \emph{prime} modulus in this family cannot exceed $\sqrt x$, landing exactly on the
BV wall, where the relevant length inequalities fail; and a \emph{composite} modulus (two prime
factors) has monodromy of type $\mathrm{SL}_2\times\mathrm{SL}_2\to\mathrm{SO}_4$, which FKMS
explicitly classify as \emph{not} gallant. \textbf{Verdict:} the theorem's hypotheses are not
met on any reachable cell; no range past $x^{1/2}$ is obtained; not T(3) progress.

\paragraph{(iv) The $r_\chi$-to-Kloosterman bilinear reduction \NEGATIVE{} (\file{erdos-straus-T-3.md}, Phases~0--4).}
Escape~2 of the three-step plan (\Cref{subsec:t3}) is carried through explicitly: the weighted
count $W(x)=\sum_p r_{\chi_1}(p+4)r_{\chi_2}(p+16)r_{\chi_3}(p+36)$ is expanded, the inner
congruence system is solved via a moving CRT residue $\alpha(d_1,d_2,d_3)$, and — after opening
Type~II via additive characters — the remainder is reduced to an explicit trilinear Kloosterman
sum, of exactly the shape treated by three 2025--26 papers: Pascadi \cite{Pascadi2026},
Milić\'evi\'c--Qin--Wu \cite{MQW2026}, and Blomer--Pascadi \cite{BlomerPascadi2026}. A short
lemma (``no bilinear cell at the wall'') shows that at the Bombieri--Vinogradov boundary itself
$\delta=0$ (writing the modulus as $x^{1/2+\delta}$) there is provably no bilinear cell at all —
one Type~II factor necessarily folds to a complete residue system. \MEASURED{} Testing the
saving ranges of all three cited theorems against the dual-interval lengths available just past
the wall: Pascadi's Corollary~1.4 length constraint is satisfied but its \emph{saving} range is
not reached until $\delta\gtrsim1/3$; Blomer--Pascadi's genuine saving over the trivial Weil
bound with balanced dual lengths first appears only at $\delta\approx0.367$
($Q\approx x^{0.87}$); MQW's theorem uses the wrong (product-argument, not AP-Poisson) kernel
shape and its length hypotheses require $Q\ge x^{4/5}$ (\file{t3\_kloosterman\_ranges.py}).
\textbf{Verdict:} no-go at the stall; the entire admissible-saving gap $x^{1/2}<Q<x^{0.87}$
remains untouched by any cited 2025--26 result; not a third escape.

\begin{remark}
The four negative results above are methodologically of a piece: each identifies a specific,
named, currently-published theorem that a naive reading might expect to bridge the
Bombieri--Vinogradov wall for this problem, and each is shown, by an explicit range or hypothesis
check rather than a vague appeal to difficulty, not to reach the cell the completion step of
T(3) needs. None of the four claims to have proved a barrier theorem; all four are consistent
with — and are read by the programme as further evidence for — the k-budget invariant of
\Cref{subsec:kbudget}, since all four available techniques supply either interval-intrinsic
input (capped) or genuinely spectral input that happens not to reach far enough in this specific
family at this specific modulus. \Cref{subsec:varyingmodulus} draws out the common structural
reason.
\end{remark}

\subsection{The varying-modulus correlation gap: a capstone diagnosis}
\label{subsec:varyingmodulus}

A later document, \texttt{erdos-\allowbreak straus-\allowbreak varying-\allowbreak modulus-\allowbreak gap.md}, steps back from the four
individual negative results of \Cref{subsec:t3-negative} and asks what they have in common. It
adds a fifth, informally recorded route to the comparison — an attempt using Maynard-type
uniform-residue sieve technology (Maynard~I--III) to reach a thin near-Bombieri--Vinogradov slab
just past the wall, which the document notes was never filed as a separate look because it fails
for the same reason as the Zheng reroute: any window inside that thin slab still forces
sacrificing one of the three kernels, which reduces to the already-closed two-modulus cell of
\Cref{subsec:t3-negative}(i). With this fifth route included, the document states its finding as
an absence rather than a construction: after Phase~1 of \texttt{erdos-\allowbreak straus-\allowbreak T-\allowbreak 3.md}
(\Cref{subsec:t3}), what remains is square-root cancellation, uniform in $d_2,d_3$, of
\[
\sum_{d_2,d_3}\chi_2(d_2)\chi_3(d_3)\,K(\cdot\,;D), \qquad D=\operatorname{lcm}(d_2,d_3)\asymp d_2d_3,
\]
where both $d_2$ and $d_3$ run up to (and, before switching, past) $\sqrt x$, so that $D$ itself
is the \emph{jointly growing, multi-factor} object — not a fixed quantity relative to which
something else varies.

\begin{claimstmt}[The two-shape classification, \CLAIMED{}]
Every method surveyed by the programme, across all five routes, falls into exactly one of two
shapes: \textbf{(1) fixed residue, growing modulus} — Bombieri--Vinogradov, BFI well-factorable
estimates, and Maynard~I--III, which all want a residue $a$ held independent of $q$ as $q$
grows, whereas the T(3) residue $\alpha(d_2,d_3)$ moves with the modulus by construction; or
\textbf{(2) fixed modulus (field), varying argument} — the $\delta$-symbol's Farey dissection
past the BV wall, and the entire $\ell$-adic trace-function lineage (Katz~\cite{Katz1988};
Fouvry--Kowalski--Michel on sums of products \cite{FouvryKowalskiMichel2015} and on trace
functions over the primes \cite{FouvryKowalskiMichel2014}; Kowalski--Michel--Sawin
\cite{KowalskiMichelSawin2017}; Fouvry--Kowalski--Michel--Sawin \cite{FKMS2026}), each of which
fixes a base field $\mathbf F_q$ (or a prime modulus $q$) and lets an auxiliary argument, tuple,
or family index vary instead.
\end{claimstmt}
No route surveyed supplies cancellation when the modulus itself is the jointly growing,
multi-factor object — this is exactly the FKMS failure mode already identified in
\Cref{subsec:t3-negative}(iii), and the document traces the same cut through the wider family
tree behind it, including Ricotta--Royer's Kloosterman paths for prime-power moduli $p^n$ with
$n\ge2$ \emph{fixed} (only $p\to\infty$ varies) \cite{RicottaRoyer2016}, and a partial, harder
adjacent case (Kowalski, uniform monodromy as both the field and a conductor grow, for
elliptic-curve twists \cite{Kowalski2006}) that the document reads as structurally closer to
what T(3) would need but not the same object. The open question is posed precisely:
\begin{quote}\itshape
Does there exist — or can one construct — a Goursat--Kolchin--Ribet-type independence statement
for a family of Kloosterman-type (or gallant) sheaves in which the modulus is the outer,
growing, multi-factor summation variable, rather than the fixed base field over which a single
sheaf's argument ranges?
\end{quote}
The document explicitly does not claim the gap is permanent — only that the specific family tree
already engaged by the programme does not contain a citation that closes it — and it notes one
untested avenue as a plausible testbed rather than a claimed result: the function-field analogue
$\mathbf F_q[t]$, where Weil-quality bounds are unconditional, so that independence across
several jointly growing polynomial moduli might be stateable and checkable before any attempt
over $\Z$. This document, more than any single one of the four negative results it synthesises,
is the programme's clearest statement of exactly where the T(3) lower bound's difficulty is
currently understood to live.

\subsection{T(A)$^+$: a uniform Selberg bound, and E\_lane}
\label{subsec:talane}

Independently of the T(3) completion question, the programme proves a genuinely unconditional
upper bound at \emph{every} fixed covering width $A$, generalising the T(3)$^+$ argument.
\begin{claimstmt}[T(A)$^+$, \CLAIMED]
For each fixed $A\ge1$, writing $\beta(A)=\sum_{a\le A}1/\varphi(4a)$ for the joint sieve
dimension (the number of distinct forbidden residues, which residue collisions can only
decrease — verified computationally to be exactly zero for $A\le200$),
\[
S(A,x) \;\ll\; C_{\mathrm{sieve}}(A)\,(\log x)^{-\beta(A)}, \qquad
C_{\mathrm{sieve}}(A) = \Gamma(\beta(A)+1)\exp\bigl(\gamma\beta(A)-B(A)\bigr),
\]
by a Selberg sieve at level $z=x^{1/4}(\log x)^{-B}$ with Bombieri--Vinogradov remainders (the
$A=3$ case is exactly T(3)$^+$).
\end{claimstmt}
A crucial supporting fact, proved as a genuine unconditional theorem (Stirling's formula applied
to the explicit bound $\beta(A)=O(\log A\log\log A)$, itself a consequence of $\varphi(n)\gg
n/\log\log n$), is that the $\Gamma$-factor's contribution to $\log C_{\mathrm{sieve}}(A)$ is
$\exp(o(\log^2A))$ — sub-critical against the $\log^2A$ covering law, so this apparatus cannot
by itself resurrect the covering density's own scaling. \MEASURED{} the Mertens-constant sum
$B(A)$ is computed at two independent cutoffs ($Q=10^6,10^7$) and agrees to three decimals,
staying in $[0.32,0.47]$ for $A\le200$; its \emph{expected} size $O(\log A)$ follows from
Norton's lemma \cite{Norton1976} together with the typical size of $L(1,\chi)$ in the
Granville--Soundararajan circle \cite{GranvilleSoundararajan2003}.
\begin{claimstmt}[Lemma FL, cited]
The dimension-uniform fundamental lemma (Halberstam--Richert \cite{HalberstamRichert1974};
Friedlander--Iwaniec \cite{OperaDeCribro}) extends T(A)$^+$'s main term and remainder uniformly
in $A$, for $A\le\exp(c\sqrt{\log x})$, with effective constants.
\end{claimstmt}
\begin{claimstmt}[E\_lane, \CLAIMED]
Optimising the Selberg exponent against Lemma~FL's growing-level constraint, there exist
effective $c,c'>0$ such that the number of hard primes $p\le x$ escaping the $d=1$
aligned-prime lane at width $A(x)=\exp(c\sqrt{\log x})$ satisfies
\[
E_{\mathrm{lane}}(x) \;\ll\; x\exp\bigl(-c'\sqrt{\log x}\,\log\log x\bigr).
\]
\end{claimstmt}
E\_lane is deliberately \emph{weaker} than Vaughan's 1970 bound: the $d=1$ lane carries only
$\log A$ worth of sieve mass, against the full covering box's $\log^2A$. It is presented as the
provable, effective ``floor'' the aligned-prime sieve alone can deliver — a useful calibration
point rather than a record.

\subsection{E\_power: a withdrawn covering-box claim}
\label{subsec:epower}

The intended strongest unconditional analytic result concerned the \emph{full}
covering box on \emph{all} integers (not only hard primes) — recall from \Cref{subsec:erdosstraus-lean}
that Lean's \lean{Covered}/\lean{Survivor} predicate on the full $(a,c,d)$ box is a
\emph{superset} of the true ES exceptional set, so an upper bound here would automatically be an upper
bound there. The claim below is \WITHDRAWN{}.

\begin{lemma}[Lemma SM, second moments, \CLAIMED{} with a written proof]
\label{lem:sm}
Let $Q(A)$ be the moduli $q=4acd-1$ of the Lean covering cells ($1\le a,c\le A$, $1\le d\le5$),
$\mu(A)=\sum_q 1/q$ the total mass, $\mu_\ell(A)=\sum_{\ell\mid q}1/q$ the mass divisible by a
prime $\ell$, $\Delta(A,T)=\sum_{\ell>T}\mu_\ell^2$, and
$\delta_{\mathrm{ub}}(A,T)=\max_q\sum_{\ell\mid q,\,\ell>T}\mu_\ell$. There is an absolute
effective $C$ such that for all $A\ge3$ and $3\le T\le20A^2$,
\[
\Delta(A,T) \;\le\; C\Bigl(\frac{\mu(A)^2}{T}+\frac{(\log2A)^2}{\sqrt T}
+\log\log(2A^2+3)-\log\log T+1\Bigr),
\]
\[
\delta_{\mathrm{ub}}(A,T) \;\le\; C\,\mu(A)\Bigl(\frac{\log2A}{T\log T}+\frac1{\sqrt T\log2A}\Bigr).
\]
\end{lemma}
The proof (reproduced in full in \texttt{erdos-\allowbreak straus-\allowbreak E-\allowbreak power.md}) is elementary: it reduces to a
divisor-in-arithmetic-progression bound $\sum_{n\le x,\,n\equiv a(q)}d(n)\ll (x/q)\log(2x)+\sqrt
x$, a weighted-harmonic corollary, a pointwise bound on $\mu_\ell$, and finally a second-moment
sum over primes $\ell>T$ split by Mertens' theorem. On the schedule $T=\mu/\eta$ for fixed small
$\eta\in(0,1/4]$, where $\mu(A)\asymp(\log A)^2$, both error terms in Suen's inequality
(\Cref{subsec:kbudget}) become $\Delta e^{2\delta}=O(\eta\mu)$ — exactly the input the density
argument needs. \MEASURED{} the implied constants in Lemma~SM are checked computationally on the
Lean covering cells at $A$ from $24$ up to $2000$ (\file{e\_power\_suen\_moments.py},
\file{e\_power\_suen\_moments\_large.py}): along the schedule $T/\mu\approx5$, the ratio
$R_\Delta=T\sum_{\ell>T}\mu_\ell^2/\mu^2$ settles at $\approx0.13$ and the Suen loss
$\Delta e^{2\delta}/\mu$ stays below $0.15$ throughout this two-order-of-magnitude range.

\begin{theorem}[E\_power, \WITHDRAWN]
\label{thm:epower}
There exist effective constants $c,\delta,X_0>0$ such that for all $x\ge X_0$ and
$A=\exp(c\sqrt{\log x})$,
\[
S_A(x,2x) \;\ll\; x^{1-\delta},
\]
where $S_A(x,2x)$ is the number of $n\in[x,2x)$ surviving the level-$A$ covering box.
Consequently the number of $n\in[x,2x)$ for which $4/n$ is \emph{not} a sum of three unit
fractions is $\ll x^{1-\delta}$.
\end{theorem}
This statement, as argued in the 21~August write-up, combined Lemma~SM's second-moment input
with Suen's inequality (giving density $u(A)\le\exp(-(1-O(\eta))\mu(A))$) and a
``Lemma Transfer'' step: at $A=\exp(c\sqrt{\log x})$, the hub modulus
$M_T=\prod_{p\le T}p$ has size $x^{\alpha/\eta+o(1)}$
(with $\alpha=\kappa_{\mathrm{cov}}c^2$), and for $c$ small enough that $\alpha/\eta<1/2$ the
main term was claimed to dominate the discrepancy over residues modulo $M_T$.
The argument is \WITHDRAWN{}: primorial conditioning does not preserve $\sum 1/q$,
$\sum_{\ell>T}\mu_\ell^2$ is not Suen's fibre $\Delta$, and $O(M_T)$ is not a
discrepancy bound. A 22~August two-stage replacement (hub exponential + fibre Suen)
failed both revival numbers through $A=32$
(\texttt{erdos-straus-E-power-decision.md}). Do not read Theorem~\ref{thm:epower}
as an improvement on Vaughan (1970). Vaughan remains the published comparison.

\subsection{Gate A: the discipline separating a written record from a formalised claim}
\label{subsec:gatea}

The programme applies an explicit editorial control, ``Gate~A'' (\texttt{erdos-\allowbreak straus-\allowbreak E-\allowbreak partial.md}),
governing what may and may not be compiled into the Lean interface
\lean{AnalyticSurvivorBound}. Gate~A \emph{admits}: the elementary Layer-A covering machinery
(already compiled); a written paper of T(A)$^+$/E\_lane/E\_power, with both ``die conditions'' of
\Cref{subsec:t_a} stated in the abstract; the Granville--Soundararajan restatement as a
dictionary into a published literature, not as extra independent characters. Gate~A
\emph{forbids}: discharging
\lean{AnalyticSurvivorBound} in Lean, or any lemma whose conclusion is zero large hard
survivors; treating E\_power or E\_partial as H\_ES or as an assault on it; retuning $\kappa$
from the flat prime-aligned Euler factor $C_{\mathrm{euler}}\approx0.91$ (\MEASURED{}, computed
in \file{c4\_euler\_factor.py} to be flat in $A$ from $A=40$ to $A=200$, with zero residue
collisions verified for $A\le200$); densifying covering while the \emph{conditional retention
probability} $\mathrm{cond}(a)=\Pr[R_a\mid R_1,\dots,R_{a-1}]$ sits at $\approx1$ (\MEASURED{} at
$x=10^9$, $\mathrm{cond}(a)\to1$ past $a=80$, mean $0.993$ on $a=161$--$200$ — meaning
essentially all of the extra covering width beyond $A\approx80$ contributes what the archive
calls \emph{dummy covering mass}, not genuine additional constraint); and running further
large-scale computation beyond $x=10^9$. This discipline is the explicit mechanism by which the
archive prevents a withdrawn covering-box claim (Theorem~\ref{thm:epower}) from being
silently over-claimed as progress toward the load-bearing interface.

\subsection{The Granville--Soundararajan reformulation}
\label{subsec:gs}

A final analytic-lane document (\texttt{erdos-\allowbreak straus-\allowbreak gs-\allowbreak reformulation.md}) restates the T(A)
survivor-density target in the language of pretentious multiplicative number theory: each
covering condition corresponds, via quadratic reciprocity, to a real Dirichlet character
condition, so the joint survival event is a large-deviation lower tail of a sum of quadratic
character values at primes — precisely the object the Granville--Soundararajan pretentiousness
programme \cite{GranvilleSoundararajan2001,GranvilleSoundararajan2007} and Lamzouri's
large-deviation theorems for character sums were built to compute. The document is explicit that
it contains \emph{no new theorem}; its role is to name the correct published dialect for the
open question of \Cref{subsec:sievedesk}, and it records two structural caveats: the number of
covering moduli in play ($A=\exp(c\sqrt{\log x})$) is small enough that a naive application of
Heath-Brown's quadratic large sieve \cite{HeathBrown1995} is trivial, so the genuine leverage
would have to come from \emph{growing moments} rather than a raw second-moment bound; and
character-sum cancellation, unlike interval-intrinsic randomness, is \emph{not} capped by the
k-budget invariant of \Cref{subsec:kbudget}, so this is one of the few investigated directions
the invariant does not pre-refute. A companion measurement, the ``C2'' cross-moduli
Jacobi-symbol correlation test (\file{c4\_c2\_symbols.py}), checks directly whether reciprocity
among the forced covering symbols produces \emph{super}-multiplicative shrinkage beyond the
deterministic exclusions already counted; \MEASURED{} at $x=10^9$, $A\le80$ it does not: mean
$|z|$ on covering-modulus pairs ($0.73$) is if anything \emph{smaller} than on inert pairs
($0.89$), and the sign of any residual bias is wrong for it to explain the observed selection
constant $\widehat C(80)=2.69$.

\subsection{The analytic worklist: an open question posed to the sieve literature}
\label{subsec:sievedesk}

The programme maintains one further document, \texttt{erdos-\allowbreak straus-\allowbreak sieve-\allowbreak desk.md}, whose own
header describes it as ``the analytic question list. It is not a theorem, not a request to prove
the Erd\H{o}s--Straus conjecture, and not a request to densify covering.'' Unlike the working
notes discussed above, it is written to be legible to a reader outside the programme: it restates
the T(A)$^+$/E\_lane/E\_power results as theorems already established in this archive, restates
the T(3) completion question and the negative findings of \Cref{subsec:t3-negative,subsec:varyingmodulus}
as the open remainder, and names the relevant published dialect as the half-dimensional/
vector-sieve literature on ker-character conditions at primes — in particular the active
$m^2+n^2+1$ literature (Hooley~\cite{Hooley1973}; Iwaniec~\cite{Iwaniec1974,Iwaniec1976};
Nath--Xie~\cite{NathXie2025}; Fuchs--Hsu--Rickards--Schindler--Stange~\cite{FHRSS2025}). It is
explicit that this is a bounded mathematical question, not a request for anyone to prove ES, and
that it should not be mixed in the same discussion with the geometric question of
\Cref{subsec:loughran-questions}: the two are addressed to different published literatures and
share only the identification, discussed in \Cref{subsec:synthesis}, that both ultimately reduce
to the same analytic core.


\section{Route 2: the Bright--Loughran geometric approach}
\label{sec:route2}

\subsection{Background: the Brauer--Manin theory of Erd\H{o}s--Straus surfaces}
\label{subsec:blbackground}

Bright and Loughran \cite{BrightLoughran2020} study, for each $n\ge2$, the affine cubic surface
$U_n:4xyz=n(xy+yz+zx)$ over $\Q$, showing $\Br U_n/\Br\Q\cong\Z/2$, generated by the
transcendental quaternion algebra $\alpha=(-u_1/u_3,\,-u_2/u_3)$, computed on the
desingularisation $\widetilde U_n$ (whose complement of three boundary lines is a torus
$\mathbb G_m^2$). Their principal results, all formalised or directly used in this archive:
\begin{itemize}[leftmargin=1.4em]
  \item \textbf{Theorem~1.2}: every natural-number solution satisfies
        $\prod_{p\mid n}(-u_1/u_3,-u_2/u_3)_p=-1$ (\Cref{subsec:brightloughran-lean});
  \item \textbf{Theorem~1.5}: non-natural integer solutions give $+1$ instead
        (\Cref{subsec:blroute-lean}, the octant invariant);
  \item \textbf{Theorem~1.8}: there is \emph{no} Brauer--Manin obstruction to ES itself — the
        conjecture survives their computation, but this is a non-obstruction statement, not an
        existence proof;
  \item \textbf{Theorem~1.9}: Brauer--Manin is \emph{not} the only obstruction to strong
        approximation on these surfaces — the exhibited failure is a consequence of the
        \emph{finiteness} of integral points (the box, Lemma~3.10) via Lang--Weil, i.e.\ a
        density phenomenon, not an obstruction-theoretic one.
\end{itemize}
The programme's own summary of the open mathematical content, after formalising this apparatus:
\emph{ES is the question of occupancy of one real component ($C_+$, the positive octant) of a
surface that already has integral points in the other} (mixed-sign points always exist,
\lean{esZ\_nonempty}).

\subsection{The formalised $\Z$-model and TUB-EP}

\Cref{subsec:blroute-lean} already describes the Lean content in detail: the affine $\Z$-model,
the box bound, the octant invariant, and the central named hypothesis
\lean{TubEpHypothesis} — occupancy of the positive octant for every $n\ge2$ — proved
\emph{equivalent} to \lean{ErdosStraus} itself (\lean{tub\_ep\_iff\_erdos\_straus}). This
equivalence is treated throughout the archive as a methodologically important negative fact
rather than a positive step: requirement R2 of the acceptance specification
(\Cref{sec:spec}) demands that a genuine load-bearing conjecture $C$ have mathematical content
\emph{beyond} ES itself, and a certified equivalence is explicitly disqualified as ``a
translation, not a target'' with ``zero strategic distance''. The geometric route's honest
status, reached only after the negative investigations below, is that it supplies
\emph{class-theory and an effectivity framework} — a precise vocabulary in which a well-posed
question about a class of surfaces can be stated — rather than a proof mechanism with tools that
presently exist.

\subsection{Harpaz's descent-fibration theorem: inapplicable, and provably so}
\label{subsec:harpaz}

The first candidate geometric mechanism tested is Harpaz's theorem on integral points of conic
log~K3 surfaces \cite{Harpaz2019}, whose engine is Swinnerton-Dyer-style descent on norm-$1$ tori
of a quadratic field, contingent on independence, in $\Q^\times/(\Q^\times)^2$, of a specific
finite collection of classes attached to the fibration. As established unconditionally in
\Cref{subsec:conicfiber} (\lean{ConicFiber.fiber\_identity}), \emph{every} conic fibre of $U_n$,
for every $n\ge2$ and every choice of fibre parameter $t$, is algebraically split: $uv=(nt)^2$.
On a split conic there is no non-trivial Selmer structure to compare, and Harpaz's hypotheses
fail degenerately rather than merely being hard to verify. The programme records this both as a
negative result for the naive application, and, via
\lean{fiber\_to\_divisor}/\lean{divisor\_to\_fiber}, as a positive structural theorem: along every
fibration, TUB-EP's existence question \emph{is} the divisor problem of \Cref{sec:covering} — the
two routes are, at the level of this one surface family, literally the same question in two
coordinate systems.

\subsection{Cao--Xu toric strong approximation: the wrong theorem}

A second natural candidate is Cao--Xu's theorem that Brauer--Manin suffices for strong
approximation on genuinely toric varieties \cite{CaoXu2018}. The desingularised surface
$\widetilde U_n$ does contain a torus $\mathbb G_m^2$ in its interior, but — as recorded in BL's
own Remark~3.12 and reconfirmed in the plan's route-2 dossier — the natural torus action does
\emph{not extend} to the boundary compactification, so $\widetilde U_n$ is not an honest toric
variety in the sense Cao--Xu's theorem requires. \LEAN{} \lean{esQ\_pos}: the existence of a
positive \emph{rational} point $(n,n,n/2)$ independently of any BM computation is recorded as a
direct illustration of why Cao--Xu cannot, by itself, be what fills the octant with an
\emph{integral} point.

\subsection{Candidate C5 and its refutation chain: the real-compatible twist}
\label{subsec:c5}

The most substantial single geometric idea investigated is candidate conjecture C5: since the
Brauer class $\alpha$ is exactly what obstructs strong approximation, perhaps descending through
$\alpha$ — replacing $U_n$ by a torsor on which $\alpha$ trivialises — produces \emph{non-split}
conic fibrations to which Harpaz's engine genuinely applies. The investigation, entirely
documented and largely kernel-checked, proceeds through three stages:
\begin{enumerate}[leftmargin=1.6em]
  \item \textbf{C5 (naive, refuted same day).} The naive $\alpha$-trivialising cover
        $w^2=-u_1u_3$ is shown, by \Cref{subsec:blroute-lean}'s octant invariant, to have
        \emph{no real points at all} over the positive octant (since $u_1u_3>0$ there),
        formalised as \LEAN{} \lean{TwistDescent.naive\_cover\_empty}. The Brauer class does not
        merely obstruct existence; it walls its own trivialising cover off from the octant
        entirely.
  \item \textbf{C5$_1$ (corrected, examined and re-scored).} The natural repair replaces the
        trivialising cover by \emph{real-compatible twists} $w^2=d\cdot u_1u_3$ ($d>0$
        squarefree), which do have real points over $C_+$ by construction. Every positive
        solution of $U_n$ is shown to lift to some twist (\lean{TwistDescent.lift}, Gate~2), with
        $d=O(n^2)$ forced by the box bound — ruling out any fixed, $n$-independent twist family
        (which would fail requirement R4, \Cref{sec:spec}). A hand calculation shows the
        pulled-back quaternion data on the twist takes the shape $(-d,-u_2u_3)$; fibering by
        $u_2=t$ and rationalising, integrality on the twist is shown to be equivalent to a
        divisor condition on norm-form values over $\Q(\sqrt{-d})$ — a genuine toolkit gain
        (imaginary-quadratic factorisation and Chebotarev machinery become available), but also,
        honestly, a restatement rather than an evasion of the same analytic core.
  \item \textbf{C5$_2$ (Gate~3, negative, closed).} The specific $t$-fibration named in C5$_1$'s
        derivation is shown to possess an \emph{explicit rational section} over $\Q(t)$
        (\Cref{subsec:twistdescent}, \lean{t\_fiber\_section\_cleared}) — the generic fibre is
        \emph{split for every $d$}, reproducing exactly the same obstruction one level up the
        tower. Harpaz's Theorem~1.0.1 (the specific, unconditional case of his theorem
        applicable here, whose hypotheses require non-split fibres and whose conclusion is in
        any case only $S$-integral with $2\in S$, not a genuine $\Z$-point of $C_+$) does not
        apply. A same-day autopsy (\lean{u1\_fiber\_section\_cleared}) shows the identical
        rational curve, merely reparametrised, is \emph{also} a section of the projections
        $\pi=u_1$ and $\pi=w$ — changing the linear projection does not open a genuinely
        non-split fibration of this surface.
\end{enumerate}
The programme records an explicit methodological lesson from this chain: \emph{run the section
test first} — before any Selmer or descent architecture is discussed for a proposed fibration of
any surface in this family, either exhibit a $\Q(t)$-section (which closes the route
immediately) or prove non-splitness of the generic fibre (which opens it). The line is
recorded as closed pending a genuinely new, non-split fibration, none having been found.

\subsection{Candidate C3 and its freeze: Campana points do not reach $C_+$}

Candidate C3 proposes a triangle-boundary semi-integral (Campana/orbifold) existence principle,
motivated by the observation that ES pairs have a totally split boundary triangle with interior
$\mathbb G_m^2$. Three checks (recorded as C3$_1$) freeze this line:
\begin{enumerate}[leftmargin=1.6em]
  \item ES pairs \emph{do} have the claimed geometric structure (a paper fact BL already use, not
        a new theorem);
  \item Campana/semi-integral points do \emph{not}, in general, imply a genuine $\Z$-point of
        $C_+$: Mitankin--Uhlemann \cite{MitankinUhlemann2025} prove that Markoff orbifold pairs
        satisfy the semi-integral Hasse principle while often having \emph{empty} integral point
        sets — their own Remark~1.5 states explicitly that Campana/Darmon points ``do not behave
        in the way that $\mathcal U_m(\Z)$ does''; consequently a class-level Campana Hasse
        principle, even with empty obstruction, does not bridge to \lean{TubEpHypothesis};
  \item there is no named theorem $T$ in the literature that supplies effective existence of a
        $\Z$-point of $C_+$ on the ES pairs by this route.
\end{enumerate}
A further claimed discriminator — that the Markoff boundary is an irreducible conic rather than a
triangle, and so C3's triangle hypothesis would automatically exclude the Markoff
counterexample — is checked directly against Mitankin--Uhlemann's actual compactification and
found to be \emph{false}: their model $x_0(x_1^2+x_2^2+x_3^2)-x_1x_2x_3=mx_0^3$ has boundary
$D=\bigcup_i\{x_0=x_i=0\}$, itself a triangle of three lines. The triangle hypothesis does not
exclude the known negative comparison case.

\subsection{The Schinzel separator programme}
\label{subsec:schinzel}

To test what an occupancy-forcing principle would need beyond mixed-sign existence and the box,
the programme studies the generalised Schinzel family $m/n=1/x+1/y+1/z$ directly.
\Cref{subsec:schinzel-lean} records the Lean content; the mathematical narrative is:
\begin{itemize}[leftmargin=1.4em]
  \item An infinite family $(2z-1)/z$ ($z\notin\{3,4,6\}$) has an explicit mixed-sign point but
        no positive solution, since positive solvability would require $(z-2)\mid4$. The
        flagship case $z=5$ is $9/5$.
  \item \MEASURED{} A bounded computer sweep of the ``narrow band'' $1<m/n<3$, $n\le40$, finds
        \textbf{86 separators} in lowest terms (mixed-sign witness exists, positive search
        provably empty — the positive search is complete for $x\le3n/m$)
        (\file{s1\_schinzel\_separators.py}); only the $9/5$ and $13/7$ instances are Lean-verified
        (\Cref{subsec:schinzel-lean}), the remaining 84 being script-level results.
  \item \textbf{The mechanism, identified precisely.} For $m/n$ bounded strictly between $1$ and
        $3$, the positive search tree has $O(1)$ branching (the leading denominator is bounded by
        $3n/m=O(1)$), so failure is a \emph{finite, Archimedean-integral} phenomenon — Bright and
        Loughran's own Theorem~1.9 finiteness mechanism, promoted from an obstruction to strong
        approximation into an outright existence failure. Ordinary Brauer--Manin machinery, which
        is not sensitive to Archimedean tree size, does not detect this.
  \item \textbf{The regime split (S1$_1$).} The Schinzel class splits into a \emph{narrow} band
        ($m/n$ bounded away from $0$), where the tree is finite and existence reduces to a
        decidable computation (\lean{decideCplus}, \Cref{subsec:schinzel-lean}), and a
        \emph{wide} band ($m/n\to0$, tree width $\sim n/2$ for $m=4$) — which is exactly where
        Erd\H{o}s--Straus itself lives. A ``wide-tree separator'' cannot, by construction, occur
        in the narrow band; extending the 86-pair sweep therefore cannot test whether any
        occupancy principle is \emph{sufficient} in the regime that matters.
  \item A hand (non-kernel) adelic sketch on $9/5$ tests whether $\alpha$ obstructs a
        $C_+$-restricted integral adele there: two explicit points ($(1,-5,1)$ at good places,
        $(9/11,45/11,3)$ as a $\Z_2$-point) are assembled with hand-computed Hilbert symbols at
        $2$, giving product of invariants $+1$ — i.e.\ \emph{no} obstruction is found on this
        sketch, meaning the $9/5$ separator's failure is invisible to the Brauer--Manin
        machinery precisely as the finite-tree mechanism above predicts. This sketch is
        explicitly labelled non-kernel and is presented as a computation to be confirmed by
        specialists, not a result.
\end{itemize}
\begin{claimstmt}[Candidate C11, refined existence schema — not proved]
\label{claim:c11}
Mixed-sign occupancy $+$ the box $+$ no integral Brauer--Manin obstruction against $C_+$ $+$
\emph{Archimedean non-degeneracy} (the admissible leading-denominator range growing with $n$,
equivalently $m/n\to0$ along the family) implies positive-octant existence.
\end{claimstmt}
C11 is scored honestly against the eight-point acceptance specification in \Cref{sec:spec} and
found to satisfy none of the eight outright; its role in the archive is as the sharpest currently
stated version of what a genuine geometric existence principle for this family would have to
assert, together with the explicit necessity proof (the 86 separators) that its fourth hypothesis
cannot be dropped.

\subsection{No Vieta: why the Markoff/Bourgain--Gamburd--Sarnak playbook cannot be imported}

\Cref{subsec:novieta} records the Lean content of the final structural finding: the ES/Schinzel
affine equation is linear, not quadratic, in each coordinate, so it admits no second root and
hence no Vieta-type involution analogous to the Markoff surface's $x\mapsto3yz-x$. This is the
formal reason the entire Bourgain--Gamburd--Sarnak strong-approximation machinery
\cite{BGS2016,GhoshSarnak2022} for Markoff-type surfaces — the natural comparison class,
since $U_n(\Z)$ is never empty exactly where sibling Markoff surfaces frequently have \emph{no}
integral points at all \cite{LoughranMitankin2021,CTWX2020} — has no orbit structure to act on
here.

\subsection{Synthesis: the two lanes are one lane}
\label{subsec:synthesis}

The geometric investigation's terminal finding (\texttt{erdos-\allowbreak straus-\allowbreak novel-\allowbreak structures-\allowbreak plan.md}
\S4v, 20 August 2026 evening) is that every geometric mechanism considered reduces to the same
analytic core:
\begin{itemize}[leftmargin=1.4em]
  \item fibrations reduce to divisor conditions (\Cref{subsec:conicfiber}), conserved under the
        real-compatible twists (\Cref{subsec:c5});
  \item the universal torsor of $\widetilde U_n$ is literally the witness equation, so
        torsor-counting (candidate C10) is the Elsholtz--Tao count \cite{ElsholtzTao2013}, whose
        all-$n$ lower bound is again the divisor-dispersion problem of \Cref{sec:route1};
  \item the affine equation's linearity (\Cref{subsec:novieta}) rules out any Vieta/Markoff
        dynamical mechanism, closing that branch of the inventory entirely.
\end{itemize}
The consequence drawn is not that geometry is worthless, but that its honest role is
\textbf{class-theory and an effectivity framework}: precise, publishable foundational questions
(\Cref{subsec:loughran-questions}) about a class of strongly-obstructed log~K3 surfaces, with an
explicit effectivity framework as a genuine open problem in its own right — not a currently
available QED engine. Both lanes converge on the same object: for every real geometric mechanism
tried, the residual open mathematics is a question about the joint distribution or divisor
structure of shifted values, exactly the object T(A)/T(3) already name analytically.

\subsection{Open geometric questions in the Bright--Loughran literature}
\label{subsec:loughran-questions}

The programme records five bounded, numbered open questions in a document,
\texttt{erdos-\allowbreak straus-\allowbreak loughran-\allowbreak orbit.md}, whose own header describes it as ``the geometric
worklist. It is not a theorem, not a request to prove the Erd\H{o}s--Straus conjecture.'' The
document identifies the relevant literature as the body of work descending from Bright and
Loughran's founding paper \cite{BrightLoughran2020} — Mitankin, Uhlemann, Colliot-Th\'el\`ene,
Wei, Xu, Harpaz, Uppal, Jahnel, and Schindler among others — and states the five questions
purely as mathematics, without addressing them to any named individual:
\begin{enumerate}[leftmargin=1.6em]
  \item whether the notation
        $\bigl(U(\mathbb A_\Z)^{C_+}\bigr)^{\mathrm{Br}} := \bigl(C_+\times\prod_p U(\Z_p)\bigr)^{\mathrm{Br}}$
        is the right formalisation of ``no integral Brauer--Manin obstruction against a real
        component'', and whether Theorem~1.8 amounts to its nonemptiness for the ES pairs;
  \item whether the two $2$-adic Hilbert symbols in the non-kernel $9/5$ adelic sketch of
        \Cref{subsec:schinzel} are correct, and whether the construction is the correct analogue
        of Theorem~1.8 on the $9/5$ Schinzel model;
  \item whether any refined obstruction in the semi-integral literature detects ``the positive
        search tree is finite and empty'' — the mechanism the Schinzel separators exhibit but
        which ordinary Brauer--Manin does not see;
  \item whether any class-level framework in this literature asserts anything in the
        \emph{wide} regime ($m/n\to0$) — the regime in which Erd\H{o}s--Straus itself lives, as
        opposed to the narrow band in which the 86 Schinzel separators were found; and
  \item whether TUB-EP, or the refined schema C11 (\Cref{claim:c11}), is a statement that
        literature would recognise, reject, or refine as a genuine class principle.
\end{enumerate}
The document is explicit about what these questions are \emph{not}: not a request to apply
Harpaz, Cao--Xu, or Theorem~1.8 (each already shown inapplicable or insufficient,
\Cref{subsec:harpaz,subsec:route2-c3}); not a request to explain why ES models have integral
points when Markoff models do not (the premise conflates two different statements,
\Cref{subsec:novieta}); and not a request to densify covering or to treat any external claimed
proof as settled (\Cref{sec:prior}). A companion table in the source document lists further,
more specialised literatures (Jahnel--Schindler and Colliot-Th\'el\`ene--Wittenberg on affine
cubics; Mitankin--Uhlemann on Campana points; the dispersion/DFI literature behind
\Cref{subsec:sievedesk}; Schinzel's own polynomial-identity literature) as adjacent questions to
be kept separate from the five above and from the analytic question of \Cref{subsec:sievedesk}.


\section{External claims and the prior archive}
\label{sec:prior}

Alongside its own investigation, the programme carries out two triage exercises against material
external to this archive: an examination of a sibling, separately maintained research effort
(referred to throughout as the ``prior Track-1 archive''), and a direct check of a 2026 preprint
claiming an elementary proof of the conjecture. Both are documented explicitly as
\emph{library/documentation, not a merge}, and neither changes the status of any result reported
in \Cref{sec:route1,sec:route2}.

\subsection{The prior Track-1 archive: what was kept, and what was not}
\label{subsec:prior-archive}

\texttt{erdos-\allowbreak straus-\allowbreak prior-\allowbreak archive.md} documents a sibling repository — not part of this
archive — pursuing a different existence-assault strategy: a ``gateway'' construction asserting
that every hard residual prime admits a bound $A\le C$ in $\mathbb F_A^\times$ (a working
constant $C=500$, with an empirical maximum $A=211$ observed through $6\times10^9$), packaged as
finite-catalogue-type (FCT) and gateway covering across hundreds of generated Lean files. The
document is explicit about the verdict: \emph{``This programme already classified that vehicle
as interval-intrinsic covering densification. Extra slices do not empty a QED-scale box''} — i.e.\
the sibling effort is read as another instance of the density-versus-universality pattern this
programme's own k-budget invariant (\Cref{subsec:kbudget}) already explains. \Cref{tab:prior}
summarises what was, and was not, brought into this archive as a result.

\begin{table}[h]
\centering\small
\begin{tabular}{@{}p{0.34\linewidth}p{0.22\linewidth}p{0.34\linewidth}@{}}
\toprule
Artefact & Kept here as & Role \\
\midrule
St\o rmer's theorem on consecutive smooth numbers (Apache~2.0; Mathlib PR~\#42040) &
  \lean{Stormer.lean} & Library theorem only; not on the T(A)/geometry path (\Cref{subsec:stormer}). \\
Bello-Hern\'andez--Benito--Fern\'andez Lemma~3/Theorem~8, and an examination of Bradford 2026 &
  \texttt{erdos-\allowbreak straus-\allowbreak bradford-\allowbreak lemma3.md} & Documentation of a finite-parameter
  covering no-go, and of why Bradford 2026 does not establish its covering claim
  (\Cref{subsec:bradford}). \\
Leochlon Type-II identity $(4b-1)(4c-1)=4p\delta+1$ &
  \texttt{erdos-\allowbreak straus-\allowbreak leochlon.md}, \lean{Leochlon.lean} & Dictionary into
  \lean{ES.IsES} (\Cref{subsec:leochlon}); not an FCT/gateway merge. \\
\bottomrule
\end{tabular}
\caption{What the prior Track-1 archive contributed to this repository.}
\label{tab:prior}
\end{table}

\noindent Explicitly \emph{not} brought into this archive, per the same document: the generated
covering cascade (\texttt{Cover*}, \texttt{Hit*}, \texttt{Band*} Lean files); the
\lean{HasGatewayNsqWitness} bounded-$A$ construction as a completeness vehicle; a
Paley--Zygmund-based ensemble argument proposed as an every-$n$ proof; a claimed
Matthiesen--Wang interface discharge of a ``ForcedSix'' obstruction (the prior archive's own
record already states that it does not discharge it); further finite-catalogue growth; and any
zkVM/SP1 verified-execution wrapping. The document's closing instruction is unambiguous: do not
treat \lean{Stormer.lean} or \lean{Leochlon.lean} as progress toward
\lean{erdos\_straus\_of\_interface}, and do not reopen the sibling archive's Bounded-$A$
construction or Bradford 2026 as a proof.

\subsection{Bradford 2026 and the Bello-Hern\'andez--Benito--Fern\'andez rigidity lemma}
\label{subsec:bradford}

\texttt{erdos-\allowbreak straus-\allowbreak bradford-\allowbreak lemma3.md} examines two related pieces of external literature
directly.

\paragraph{Lemma~3 / Theorem~8 (Bello-Hern\'andez--Benito--Fern\'andez).} The statement, quoted
from \cite{BelloHernandez2026} Theorem~8 with the proof deferred to \cite{BelloHernandez2010}
Lemma~3, is a rigidity
result: one cannot generate a finite number of congruence classes containing \emph{all} primes
$p\equiv1\pmod4$ by fixing, in the classical Type~I/Type~II witness equations, three of the four
parameters to a finite set and leaving the fourth free. The proof idea is elementary lcm
arithmetic — each frozen triple produces an arithmetic progression or finite set of $n$; taking
the lcm $T$ of the moduli attached to any finite collection $S$ of triples, the progression
$\{4Tt+1:t>0\}$ cannot arise from any triple in $S$, since matching the parametric formula would
force $1$ itself to be an Erd\H{o}s--Straus number. The document records precisely what this does
and does not rule out: it rules out any \emph{finite} catalogue of fixed-parameter identities,
matching this programme's own ``do not densify covering'' discipline exactly; it does
\emph{not} rule out constructions whose free parameters grow with $n$ (there being then no
single finite $S$ and no single modulus $T$) — such constructions merely lie outside the
lemma's hypothesis, which is not the same as a confirmed density escape.

\begin{claimstmt}[Bradford 2026, \NEGATIVE]
K.~Bradford's preprint \cite{Bradford2026} claims an elementary proof of ES for every prime,
treating $p\equiv1\pmod4$ via Type~I/Type~II identities and asserting that combining them yields
a covering system. The programme's examination finds the covering claim unsubstantiated: the
preprint's only written verification is a $k=0$ list of four moduli ($44,20,8,140$) and the
observation that $5,13,17,29$ are the first four primes $\equiv1\pmod4$; there is no verification
that every such prime lies in one of the listed progressions, no table for any $k\ge1$, and no
engagement with Mordell's quadratic-residue obstruction. Three specific reasons the covering
claim fails are recorded: (i)~Mordell's classical obstruction (\Cref{subsec:whatitis}) already
forbids any finite family of polynomial identities from covering $p\equiv1\pmod{840}$, since $1$
is always a square, and Bradford's Lemmas~1--2 are identities of exactly this type;
(ii)~the written proof exhibits only $k=0$, so it is precisely the finite-parameter cover that
Lemma~3 above rules out, and no infinite template over all $k$ is supplied;
(iii)~listing four small primes is not a covering argument; and (iv)~several of the listed
arithmetic progressions ($13\bmod20$, $29\bmod44$) coincide with identities already present in
the classical Mordell layer of this archive's own Lean development
(\Cref{subsec:mordell}).
\end{claimstmt}
The document's conclusion is that formalising Bradford's claim as a proof is not worthwhile, and
that formalising its Lemmas~1--2 as general identity schemas would itself constitute the
covering densification this programme deliberately avoids; consequently neither is carried into
Lean. Three further Bradford-adjacent publications \cite{Bradford2021,BradfordIonascu2015,Bradford2024}
are noted as distinct papers that do not remove the
Mordell obstruction either.


\section{The candidate-conjecture catalogue}
\label{sec:candidates}

\texttt{erdos-\allowbreak straus-\allowbreak candidate-\allowbreak conjectures.md} maintains a running list of named candidates for
the load-bearing statement $C$ that a proof of ES would ultimately rest on, each scored against
the acceptance specification of \Cref{sec:spec} and each carrying an explicit statement of its
\emph{primary risk} — the requirement most likely to kill it. The document's own framing is
important: ``None is asserted true; several are expected to fail. That is what a candidate list
is for.'' \Cref{tab:candidates} summarises the full list and its current status; candidates
discussed in detail elsewhere in this report are cross-referenced rather than re-derived.

\begingroup
\small
\begin{longtable}{@{}p{0.08\linewidth}p{0.27\linewidth}p{0.17\linewidth}p{0.38\linewidth}@{}}
\toprule
\textbf{ID} & \textbf{Statement (informal)} & \textbf{Input class} & \textbf{Status / primary risk} \\
\midrule
\endhead
C1 & Structured-modulus dispersion: level of distribution~1 for the joint divisor condition,
     uniformly over the growing $(a,d)$ family & spectral/bilinear
   & Reference candidate; ``H\_ES in its most attackable dress''. Risk R2: may be ES relocated. \\
C2 & Reciprocity coupling: forced character conditions compound \emph{beyond} their raw covering
     mass via quadratic reciprocity & algebraic rigidity
   & \MEASURED{} \textbf{Does not fire} (\Cref{subsec:gs}): cross-moduli Jacobi-symbol
     correlations are noise, wrong sign for $\widehat C>1$. \\
C3 & Triangle-boundary semi-integral (Campana) existence principle & global geometry
   & \textbf{Frozen at C3$_1$} (\Cref{subsec:route2-c3}): Campana points do not reach $C_+$; the
     claimed Markoff-excluding discriminator is false. \\
C4 / C4$_1$ & Simultaneous shift-roughness / \textbf{T(A)}: joint density of class-roughness at
     fixed covering width & spectral (Wirsing/LSD) + compositum densities
   & \textbf{The programme's primary in-house line} (\Cref{subsec:t_a}); T(3) opening case
     claimed at the upper-bound half (\Cref{subsec:t3}). \\
C5 / C5$_1$ / C5$_2$ & Descent through $\alpha$: real-compatible twist $V_{n,d}$ & global
     geometry (Harpaz-type)
   & \textbf{Refuted / closed} (\Cref{subsec:c5}): the named $t$-fibration is split
     (Gate~3, kernel-checked); ``run the section test first''. \\
C6 & Theory-4 flatness: Selberg-$L^2$ residual of the survivor indicator is $\varepsilon_j$-flat
     against later covering shells & spectral (new equidistribution)
   & Rungs~1--3b partially executed (\Cref{subsec:kbudget} discussion); primary risk R2/R3, may
     be H\_ES relocated at asymptotic scale. \\
C7 / C7$_1$ & Bounded-coupling covering family: a new identity family breaking the Poisson
     coupling law & none (combinatorial discovery)
   & \MEASURED{} \textbf{First search negative} (\Cref{subsec:computation-misc}): 193 degree-1
     identities searched, none covers a hard class; bounded-$\omega$ mass does not break the
     coupling law at toy scale. Probably false. \\
C8 & Thin-exception mop-up: a two-stage Frobenius-thin containment plus effective Chebotarev
     closure & spectral + rigidity
   & Not pursued computationally; primary risk is the classical ``for-all gap'' at R4 (thinness
     historically gives density zero, not emptiness). \\
C9 & Discriminant-uniform spectral equidistribution (theta route): a $p$-dependent uniform
     DFI-type statement & spectral \newline (GL(2)/\allowbreak metaplectic)
   & Not pursued; primary risk is that $p$-dependent uniformity at this level is beyond current
     spectral technology. \\
C10 & Effective uniform log-Manin on the ES Cox ring (universal torsor counting) & mixed
     geometry/circle method
   & \textbf{Identified as C1/C4 in disguise} (\Cref{subsec:synthesis}): the torsor is the
     witness equation, so this is the Elsholtz--Tao count. \\
C11 & TUB-EP-refined: mixed-sign $+$ box $+$ no integral BM obstruction against $C_+$ $+$
     Archimedean non-degeneracy $\Rightarrow$ positive existence & global geometry (schema)
   & \Cref{claim:c11}; the sharpest current geometric schema, satisfies none of R1--R8 outright;
     necessity of the fourth hypothesis proved via the 86 Schinzel separators. \\
NV & No Vieta (closes the geometric mechanism inventory) & --
   & \LEAN{} \textbf{Proved} (\Cref{subsec:novieta}): the affine equation is linear in each
     coordinate; no second root, no Markoff-type involution. \\
S1 / S1$_1$ & Schinzel separator suite; narrow/wide regime split & --
   & \Cref{subsec:schinzel}: 86 separators found; the wide band is where ES itself lives; a
     non-kernel adelic sketch on $9/5$. \\
\bottomrule
\end{longtable}
\endgroup
\label{tab:candidates}

\subsection{C3, restated for completeness}
\label{subsec:route2-c3}
See \Cref{sec:route2} for the full discussion of C3/C3$_1$, C5/C5$_1$/C5$_2$, S1/S1$_1$, C11 and
NV, and \Cref{sec:route1} for C1/C2/C4/C4$_1$/C6/C10. The programme's own portfolio note groups
C1, C4, C4$_1$, C6, C9 and C10 as ``one object'': they die together if the selection constant
$C(A)$ turns out to be super-critical in $\log^2A$ ($c'\ge\kappa$), or if the conditional
retention probability $\mathrm{cond}(a)$ sticks at~$1$ (dummy covering) out to QED scale; as
measured (\Cref{subsec:t_a}, \Cref{subsec:gatea}), growing $A$ to $200$ at $x=10^9$ showed
neither has yet happened, but the margin against $\kappa$ is thin ($c'=0.104$ against
$\kappa=0.139$) and the dummy-covering trend is the more immediate concern.


\section{The acceptance specification}
\label{sec:spec}

\texttt{erdos-\allowbreak straus-\allowbreak conjecture-\allowbreak spec.md} formulates, as an explicit deliverable in its own right,
an \emph{acceptance specification} for what any candidate load-bearing conjecture $C$ must
satisfy in order to genuinely solve ES rather than merely restate it. The document frames this
as itself a piece of mathematics: ``Each requirement is traceable to a theorem, measurement, or
refutation in the programme record; each has already eliminated at least one proposed
candidate.'' The eight requirements, with the notation $A$ for the covering level, $\kappa\approx
0.139$ for the measured marginal constant, and $q\mid p+4a^2d,\,q\equiv-1\pmod{4ad}$ for the
divisor form:

\begin{description}[leftmargin=1.6em]
  \item[R1 (Verified sufficiency).] $C$ must imply \lean{HardLandingHypothesis} — hence ES —
        through kernel-checked implications. \emph{Status: cheap and fully under the
        programme's control}; both named interfaces satisfy it by construction
        (\Cref{subsec:coveringlayer}).
  \item[R2 (Independent content).] $C$ must be attackable by methods that do not already amount
        to proving ES; a certified equivalence is a translation, not a target.
        \emph{Status: the open risk for essentially every candidate} — \lean{TubEpHypothesis}
        fails it outright (\Cref{subsec:blroute-lean}); T(A)/C1's relocation risk is flagged as
        needing an early expert check.
  \item[R3 (Barrier evasion — the k-budget invariant).] $C$'s proof must inject arithmetic input
        from \emph{outside} the interval-intrinsic class capped by \Cref{subsec:kbudget}:
        admissible classes are spectral/bilinear arithmetic, algebraic rigidity used
        non-statistically, or genuinely global geometry not routed through fibres. \emph{Any $C$
        whose natural proof is ``sieve harder'' fails R3 before it starts.}
  \item[R4 (Exceptional-set annihilation, with mandatory growth).] Fixed parameter boxes
        provably leave $\sim x/(\log x)^{1+c(D_0)}$ exceptions (compounding failure densities
        across a fixed box); $C$ must involve a level growing with $x$ at rate
        $\log^2(\mathrm{level})\gtrsim\log x$, i.e.\ $\mathrm{level}=\exp(c\sqrt{\log x})$ with
        $\kappa c^2>1$. \emph{Status: eliminates every fixed-box proposal outright}, including an
        externally advised fixed-$D_0$ landing lemma (shown heuristically false,
        \Cref{subsec:kbudget} discussion of plan \S4q).
  \item[R5 (Splitness compatibility).] $C$ must locate the arithmetic in integrality and
        divisors, not in fibre-level local--global structure, since every conic fibre of the ES
        surface is split (\Cref{subsec:conicfiber}). \emph{Status: satisfied natively} by the
        divisor form (\Cref{sec:covering}); \textbf{fails} for every fibration-based geometric
        mechanism tried (\Cref{subsec:c5}).
  \item[R6 (Consistency with known negatives).] $C$ must not prove too much: it must be
        compatible with genuine Markoff-type integral Hasse failures
        \cite{GhoshSarnak2022,LoughranMitankin2021,CTWX2020}, with BL Theorem~1.9, and with
        abelian insufficiency (no finite set of congruence conditions covers the hard classes).
        \emph{Status:} eliminates C3 as originally stated (\Cref{subsec:route2-c3}).
  \item[R7 (Effectivity).] $C$ must come with an explicit threshold $X_0$ — an ineffective
        Siegel-style existence statement would establish ES as true while leaving the
        machine-checked proof permanently uncompilable. \emph{Status: satisfied by the finite
        certificate mechanism of \Cref{subsec:erdosstraus-lean} in principle; open for every
        geometric candidate.}
  \item[R8 (Statability now, in a published dialect).] $C$ must be formalisable as
        a kernel \lean{Prop} today, and phrased in a published dialect — bilinear forms in
        divisor problems, or global existence principles for log~K3 surfaces — so that the
        statement has existing theorems to cite, not only ES. \emph{Status: both named
        interfaces satisfy (a); the analytic worklist (\Cref{subsec:sievedesk}) and the
        geometric questions (\Cref{subsec:loughran-questions}) are the programme's record of
        (b).}
\end{description}

\Cref{tab:scorecard} reproduces the specification document's own scorecard for the two live
candidates it regards as closest to satisfying all eight requirements simultaneously.

\begin{table}[h]
\centering
\small
\begin{tabular}{@{}p{0.20\linewidth}p{0.36\linewidth}p{0.36\linewidth}@{}}
\toprule
Requirement & \lean{DivisorLandingBound} (growing bound) & Existence principle $\Rightarrow$ TUB-EP \\
\midrule
R1 sufficiency & \textbf{proved} & via \lean{erdos\_straus\_of\_tub\_ep} once $C$ is stated \\
R2 independence & open --- the known risk & open --- if stated for a genuine surface class \\
R3 barrier evasion & designed for spectral input; unproven & automatic (not interval-bound) \\
R4 growth & \textbf{built in} & must be engineered (uniformity in $n$) \\
R5 splitness & \textbf{native} (it is the divisor form) & the hard design constraint \\
R6 negatives & quantitative, not local--global & \textbf{the hard design constraint} \\
R7 effectivity & plausible (dispersion proofs usually effective) & needs care \\
R8 statability/dialect & \textbf{stated; dispersion school} & statable schema; Bright--Loughran literature \\
\bottomrule
\end{tabular}
\caption{The acceptance-specification scorecard, reproduced from \texttt{erdos-\allowbreak straus-\allowbreak conjecture-\allowbreak spec.md}.
Neither candidate currently satisfies all eight requirements.}
\label{tab:scorecard}
\end{table}

The specification's own one-paragraph summary captures the programme's central methodological
stance precisely enough to reproduce verbatim: \emph{``A conjecture that solves Erd\H{o}s--Straus
must be a kernel-statable proposition with mathematical life beyond ES, whose proof injects
spectral, rigidity, or global-geometric input outside the interval-intrinsic class, which
annihilates (not merely thins) the exceptional set through a parameter growing like
$\exp(c\sqrt{\log x})$ with $\kappa c^2>1$, which respects the surface's total splitness by
living in divisor/integrality structure or in genuinely global geometry, which is consistent
with the Markoff failures, Theorem~1.9, and abelian insufficiency, which names its threshold,
and which is written in a published dialect with existing theorems, not only ES. Everything short
of this is infrastructure — valuable, publishable, and already substantially built — but not the
keystone.''} The specification document's own closing attribution, which this report reproduces
in place of an earlier, superseded version naming specific model systems, reads: \emph{``Prepared
as part of the Riley Betts Ltd Erd\H{o}s--Straus programme, facilitated by Martyn Riley
(automation and high-level strategy; not a claim of expertise in number theory, combinatorics,
or geometry). The mathematical work is a collaboration of Anthropic agents, Grok~4.6, and
Cursor agents, with Lean~4, Mathlib, and numerical test resources. No proof of the conjecture is
claimed. Every claim above is traceable to a kernel-checked theorem, a labelled measurement, or a
recorded refutation in the programme archive.''}


\section{Computational investigations}
\label{sec:computation}

Every quantitative claim in \Cref{sec:route1,sec:route2} is backed by one of twenty Python
scripts in the archive, each producing a JSON artefact recording its parameters and results. This
section describes each script's purpose and principal findings; \Cref{app:files} additionally
lists every script and its output file(s) for reference. All computation is capped, by explicit
programme discipline (``Do not run $x=10^{10}$''), at $x\le10^9$; this ceiling is stated
repeatedly and is treated as a hard methodological boundary rather than a resource constraint —
the archive is explicit that further scans at the same scale would refine constants without
resolving the open die-conditions.

\subsection{The T(A)/ClassRough family}

Twelve scripts instrument the T(A) investigation of \Cref{subsec:t_a}.

\begin{description}[leftmargin=1.6em]
  \item[\file{c4\_S\_xscan.py}] The primary $x$-scan of $S(A,x)$ at fixed $A\in\{40,80\}$ across
        $x=10^5$ to $10^9$ ($1{,}587{,}581$ hard primes at the top scale). Provides the sieve of
        hard primes and the \lean{ClassRough}-equivalent Boolean predicate reused by most later
        scripts. \MEASURED{}: $\log S$ is linear in $\ln\ln x$ ($R^2=0.997$ for $x\ge10^6$); the
        selection constant climbs $\widehat C(80):1.57\to2.23\to2.69$ across $x=10^6,10^8,10^9$.
  \item[\file{c4\_growing\_A.py}] Growing-$A$ scan at fixed $x=10^9$, $A\le200$. \MEASURED{}: the
        two central die-condition quantities of \Cref{subsec:gatea} —
        $\log\widehat C(A)\sim0.104\log^2A$ on $A\ge40$ ($R^2=0.998$), and $\mathrm{cond}(a)\to1$
        past $a=80$ (mean $0.993$ on $a=161$--$200$).
  \item[\file{c4\_growing\_box\_probe.py}] An earlier, coarser probe of the same
        quasi-independence question in both the $d=1$-growing-$A$ and $2$-dimensional
        $(a,d)$-growing-box directions, predating and superseded in detail by the two scripts
        above.
  \item[\file{c4\_two\_shift\_probe.py}] Pairwise independence: $\rho_{ab}/(\rho_a\rho_b)$ for
        $1\le a<b\le A\le120$. \MEASURED{}: consistently $1.000\pm0.05$ — the empirical basis for
        the programme's repeated methodological point that pair-level statistics are the
        ``wrong moment'' for T(A), since the full $A$-fold intersection can drift while all
        pairs stay independent.
  \item[\file{c4\_euler\_factor.py}] The prime-aligned local Euler factor $C_{\mathrm{euler}}(A)$
        of \Cref{subsec:gatea}, computed over primes $q\le10^7$. \MEASURED{}: flat at
        $\approx0.91$ (the $1/q$-normalised column) from $A=40$ to $A=200$, with an exhaustive
        residue-collision check finding \emph{zero} collisions for $A\le200$.
  \item[\file{c4\_composite\_layer.py}] The genuine-extra composite Euler product
        $\Pi_{\mathrm{gen}}(a,\sqrt x)$ (composite aligned divisors with no aligned prime
        factor), isolating the ``$2.7$--$3.0\times$ totient mismatch'' between the prime-aligned
        and full \lean{ClassRough} marginals. \MEASURED{}: matches the empirical ratio
        $\rho_a^{\mathrm{CR}}/\rho_a^{\mathrm{prime}}$ to within $3\%$ at $A=40$ across
        $x=10^6,10^8,10^9$.
  \item[\file{c4\_certificates.py}] The certificate-converse and shared-certificate ``dummy
        covering'' tests of \Cref{subsec:t_a}. \MEASURED{}: converse holds for only $84.6\%$ of
        survivors at $x=10^6$, $A\le40$ (exact at $\omega=1$, falling to $35\%$ at $\omega=4$);
        certificates inherited from divisors of $a$ cover only $\approx73\%$ of later survivors,
        ruling out shared-certificate inheritance as the explanation for
        $\mathrm{cond}(a)\to1$.
  \item[\file{c4\_defect\_layer.py}] Exact enumeration (or Monte Carlo beyond a size cutoff) of
        the combinatorial defect between $\mathrm{ClassRough}$ and the certificate union,
        confirming that the two coincide exactly on $a\in\{1,2,3,6\}$ (the elementary
        $2$-group locus $\lambda(4a)\mid2$) and diverge starting at $a=4$, matching
        \Cref{subsec:t_a}.
  \item[\file{c4\_c2\_symbols.py}] The C2 cross-moduli Jacobi-symbol correlation test of
        \Cref{subsec:gs}, at $x=10^9$, $A\le80$, $564$ survivors. \MEASURED{}: correlations
        consistent with independence; wrong sign to explain $\widehat C>1$.
  \item[\file{c4\_surface\_fit.py}] Fits the two-dimensional surface $\widehat C(A,x)$ jointly in
        $A$ and $x$ from the JSON outputs of the scripts above, isolating the $x$-slope deficit
        ($\kappa_{\mathrm{eff}}\approx0.121$ at $A=80$) from the $A$-curvature
        ($\kappa_{\mathrm{eff}}\approx0.035$ at $x=10^9$) as genuinely different partial
        derivatives of $\log S$, and tracking $c'(x)$ across $x=10^6,10^8,10^9$
        ($0.063\to0.076\to0.094$, still rising, with a naive linear extrapolation to $\kappa$
        landing around $x\sim10^{24}$ — explicitly flagged as \emph{not} a reason to run further
        large-scale computation).
  \item[\file{c4\_sieve\_constant.py}] Computes the Selberg constant $C_{\mathrm{sieve}}(A)$ and
        its Mertens-constant component $B(A)$ of T(A)$^+$ (\Cref{subsec:talane}), cross-checked
        at cutoffs $Q=10^6$ and $Q=10^7$.
  \item[\file{c4\_t3\_euler.py}] The four certificate-assignment Euler products specific to T(3)
        (\Cref{subsec:t3}), each computed to $q\le10^6$, together with the empirical occupancy
        of the four assignments among $2{,}370$ hard primes at $x=10^6$ (exact match to the
        kernel theorem's prediction: 656 class-rough on all three slices, 656 in the certificate
        union).
\end{description}

\subsection{The E\_power second-moment family}
\label{subsec:computation-epower}
\begin{description}[leftmargin=1.6em]
  \item[\file{e\_power\_suen\_moments.py}] Checks Lemma~SM's implied ratios $R_\Delta,R_\delta$
        on the Lean covering cells at $A\in\{24,48,80,120,160,200,240\}$ against several hub
        cutoffs $T$; establishes the schedule $T\approx5\mu$ used in the final theorem.
  \item[\file{e\_power\_suen\_moments\_large.py}] A two-pass, memory-efficient re-implementation
        extending the same check to $A\in\{400,800,1200,1600,2000\}$, confirming that
        $R_\Delta$ settles near $0.13$ and the Suen loss stays under $0.15\mu$ throughout a
        further order of magnitude in $A$ (\Cref{subsec:epower}).
\end{description}

\subsection{The T(3)-completion range-check family}
Four scripts, all deliberately \emph{range checks} rather than covering runs, support the four
negative results of \Cref{subsec:t3-negative}:
\begin{description}[leftmargin=1.6em]
  \item[\file{t2\_zheng\_ranges.py}] Evaluates Zheng's piecewise exponent functions
        $\mathcal L(\theta)$ on a $181\times181$ grid of the switched divisor-exponent plane,
        computing the fraction of the plane covered by ordinary Bombieri--Vinogradov, by Zheng's
        Theorem~1.1, and by neither; also implements the ``uneven leftover'' reroute bookkeeping
        of \Cref{subsec:t3-negative}(i).
  \item[\file{t3\_delta\_ranges.py}] Implements the Farey--$k$ incompatibility lemma of
        \Cref{subsec:t3-negative}(ii) and checks it symbolically across the three named cells
        (BV wall, uneven, symmetric).
  \item[\file{t3\_kloosterman\_ranges.py}] Computes the admissible-saving windows of Pascadi,
        Blomer--Pascadi and MQW as functions of the stall parameter $\delta$, locating the first
        genuine saving points cited in \Cref{subsec:t3-negative}(iv).
  \item[\file{t3\_monodromy\_ranges.py}] Checks the FKMS length hypotheses (Theorems~1.3/1.4)
        against the three named cells, confirming the mismatch tabulated in
        \Cref{subsec:t3-negative}(iii): the one cell where the inequalities hold cannot carry a
        prime modulus after switching, and the one cell that can has Type~II lengths too long
        for either theorem.
\end{description}

\subsection{Geometric and combinatorial searches}
\label{subsec:computation-misc}
\begin{description}[leftmargin=1.6em]
  \item[\file{s1\_schinzel\_separators.py}] The Schinzel-separator sweep of
        \Cref{subsec:schinzel}: exhaustive complete search for positive solutions
        ($x\le3n/m$) and a bounded search (coordinates $\le60$) for mixed-sign witnesses, over
        all reduced pairs $1<m/n<3$, $n\le40$; reproduces the 86-separator count.
  \item[\file{c7\_bounded\_coupling\_search.py}] Two bounded computations for candidate C7
        (\Cref{sec:candidates}): an exhaustive search over degree-$1$ Witness identities with
        modulus $\le30$ and coefficients $\le4$ (193 identities found, none covering a hard
        class), and a coupling-mass comparison between the full Type-I covering family and its
        bounded-$\omega(q)$ sub-slice at $A\le64$, finding no evidence of the mass
        concentration C7 would need.
\end{description}

\subsection{Methodological remarks}
Three features of the computational programme deserve note as part of its overall
epistemic discipline. First, every script's docstring states explicitly what the computation
does \emph{not} establish — e.g.\ \file{c4\_S\_xscan.py}'s header notes it neither proves nor
refutes H\_ES, and \file{s1\_schinzel\_separators.py}'s header notes it does not check
Brauer--Manin pairings. Second, several scripts are explicitly framed as adjudicating between
two or more competing predictors for the same quantity (e.g.\ \file{c4\_surface\_fit.py}
comparing $\log^2A$, $\sum\log n_\chi/a$, $\log Z_{\mathrm{global}}$, and $\sum\log n_\chi$ as
candidate formulae for $\log\widehat C(A)$, finding $\log^2A$ the best fit by $R^2$ while
explicitly declining to treat this as more than curve-fitting). Third, the repeated instruction
``do not run $x=10^{10}$'' recurs verbatim across a dozen documents and script docstrings,
functioning as an explicit acknowledgement that the open die-conditions of
\Cref{subsec:t_a,subsec:gatea} cannot be resolved by further computation at reachable scale, and
that the correct next step is theoretical (Euler-product-level) rather than empirical.


\section{Discussion and assessment}
\label{sec:discussion}

\subsection{What has been achieved}

Judged against its own stated objectives (\Cref{sec:intro}), the programme delivers five
concrete, independently assessable items. First, a genuine machine-checked reduction: the
Erd\H{o}s--Straus conjecture is formally reduced, inside the Lean kernel and with an explicit,
published axiom audit, to a single named analytic proposition
(\lean{AnalyticSurvivorBound}) together with a finite verified range now standing at
$n<100{,}000$. This is not a new mathematical fact — the reduction to the six hard classes mod
$840$ is classical — but its complete, axiom-tracked formalisation, including the covering
system and the divisor-form duality, is a genuine piece of infrastructure not previously
available. Second, a recorded negative at the covering box: Theorem~\ref{thm:epower} was
claimed as a power saving $O(x^{1-\delta})$ improving Vaughan (1970), and is \WITHDRAWN{}
(one-stage CDL and the two-stage hub replacement both fail; see the dated notice).
Third, a substantial and unusually disciplined formalisation of the Bright--Loughran
Brauer--Manin theory, including a from-scratch finite verification of the arithmetic heart of
their Lemma~3.8 and an extensive discharge of the Hilbert reciprocity the ES problem actually
needs. Fourth, and arguably the programme's most distinctive contribution, a coherent explanation
— the k-budget invariant — of \emph{why} the classical approaches have stalled, instantiated
concretely five separate times (the fundamental lemma, the large sieve, Vaughan's exponential
sums, a circuit-complexity fooling argument, and the four-to-five T(3)-completion look
investigations), each converging on the same obstruction from a different technical direction,
with \Cref{subsec:varyingmodulus}'s capstone document giving the clearest single statement of
why. Fifth, a disciplined external triage (\Cref{sec:prior}): a sibling covering-density effort
is checked against the programme's own k-budget diagnosis and found to be another instance of it
rather than an escape, and an external claimed proof is checked directly against the classical
Mordell obstruction and shown not to establish the covering statement it needs — in both cases
with the specific point of failure named rather than the claim being dismissed by assertion.

\subsection{What has not been achieved, stated plainly}

The programme is explicit, at every level of its own documentation, that it has not proved the
conjecture, has not reduced it to a statement the mathematical community regards as tractable,
and has not discharged either of its two named open interfaces. Three specific limitations are
worth isolating for a reader assessing the work's significance. First, the k-budget invariant
itself is explicitly labelled ``conjectural, proof-shaped'' rather than a theorem; it is a
strong organising heuristic, well corroborated by five independent instantiations, but it is not
established that \emph{every} conceivable interval-intrinsic method is capped by it. Second, the
T(A)/T(3) selection-constant question — whether $C(A)$ grows sub-critically against
$\exp(\kappa\log^2A)$ — remains open with a thin empirical margin ($c'=0.104$ against
$\kappa=0.139$ at the largest tested scale) and an explicitly acknowledged live risk (``dummy
covering'', $\mathrm{cond}(a)\to1$) that the programme itself regards as the more pressing
concern. Third, the geometric route's terminal finding — that every mechanism tried reduces to
the analytic core — while a genuine and carefully earned negative result, means Route~2 does not
currently supply an independent path to the conjecture; its value is prospective (an
effectivity framework for strongly-obstructed log~K3 surfaces) rather than realised.

\subsection{On the methodology}

Two methodological features of the archive merit comment independent of the mathematical
content. First, the discipline of recording refuted intermediate claims rather than silently
correcting them (most visibly in the \S4f/\S4h correction of the original, erroneous
Vaughan-exponent ``barrier'', and in the C5/C5$_1$/C5$_2$ refutation chain) makes the research
log unusually auditable: a reader can verify not only what is currently claimed but what
specific error the current claim corrects, and why. Second, the explicit ``Gate A'' discipline
(\Cref{subsec:gatea}) separating a written, useful record (E\_power) from a machine-checked claim
against the load-bearing interface is a genuine safeguard against the most common failure mode
of ambitious computational mathematics programmes — the silent inflation of a genuine partial
result into an implicit claim of having solved the real problem. Both practices are followed
consistently throughout the archive and are reflected, where relevant, in the status tags used
throughout this report.

\subsection{Summary}

In summary: the archive contains no proof of the Erd\H{o}s--Straus conjecture and does not claim
one. It contains a substantial, carefully audited Lean~4 formalisation of the classical
reduction and of a non-trivial slice of the Bright--Loughran geometric theory; a recorded
negative at the covering box (Theorem~\ref{thm:epower} is \WITHDRAWN{}; Vaughan (1970) remains
the published comparison); a precisely
argued (if not fully proved) explanation of the structural obstruction the conjecture presents to
every method tried; and four honestly negative results checking specific, currently published
2025--26 analytic number theory against that obstruction. Considered as a piece of research
infrastructure and a research log, it is unusually complete and unusually disciplined about the
distinction between what is proved, what is claimed, what is measured, and what has been tried
and found not to work.


\appendix
\section{Complete archive file catalogue}
\label{app:files}

This appendix lists every file in \file{erdos-straus-main} (87 files), grouped by type, with a
one- or two-line description and a cross-reference to the section of this report discussing it.
File sizes are as extracted from the archive.

\subsection{Lean source files (13)}
\begingroup\small
\begin{longtable}{@{}p{0.28\linewidth}p{0.08\linewidth}p{0.50\linewidth}p{0.08\linewidth}@{}}
\toprule
\textbf{File} & \textbf{Lines} & \textbf{Contents} & \textbf{\S} \\
\midrule
\endhead
\file{ErdosStraus.lean} & 3420 & Core reduction, covering system, analytic interfaces,
  \lean{ClassRough}/certificate theory, elementary $d=1$ Henriot-interface sieve. & \ref{subsec:erdosstraus-lean} \\
\file{BrightLoughran.lean} & 407 & Serre 2-adic symbol data, finite heart of BL Lemma~3.8,
  \lean{InvariantData} interface, Theorem~1.2 anatomy, $p=1009$ instances. & \ref{subsec:brightloughran-lean} \\
\file{ErdosStrausQR.lean} & 170 & Mathlib discharge: \lean{IsPrime}$\leftrightarrow$\lean{Nat.Prime},
  Yamamoto fingerprint, quadratic reciprocity. & \ref{subsec:erdosstrausqr-lean} \\
\file{ErdosStrausBLRoute.lean} & 3675 & $\Z$-model, box bound, octant invariant,
  \lean{TubEpHypothesis}, extensive Hilbert reciprocity discharge. & \ref{subsec:blroute-lean} \\
\file{ConicFiber.lean} & 194 & Split-fibre identity and exact divisor correspondence, all $n\ge2$. & \ref{subsec:conicfiber} \\
\file{TwistDescent.lean} & 219 & Real-compatible twist $V_{n,d}$; Gate~2/Gate~3; the C5$_2$ refutation. & \ref{subsec:twistdescent} \\
\file{SchinzelSep.lean} & 84 & Signed $9/5$ point; ordered no-positive-solution proof. & \ref{subsec:schinzel-lean} \\
\file{SchinzelDecide.lean} & 268 & Decidable positive-octant solver with soundness/completeness; $9/5$, $13/7$, $4/5$ instances. & \ref{subsec:schinzel-lean} \\
\file{NoVieta.lean} & 156 & Linearity of the affine equation; uniqueness; no Vieta involution. & \ref{subsec:novieta} \\
\file{EPower.lean} & 3323 & Combinatorial covering core (hub $q=s\ell$, fibre $1/\ell$, pair mass, Janson I counting). H1/H2 uninhabited; two-stage density a recorded negative. & \ref{subsec:epower} \\
\lean{Leochlon.lean} & 95 & Type-II identity dictionary $(4b-1)(4c-1)=4p\delta+1\to$\lean{ES.IsES}. & \ref{subsec:leochlon} \\
\lean{Stormer.lean} & 855 & St\o rmer's theorem on consecutive smooth numbers (ported, Apache~2.0, library only). & \ref{subsec:stormer} \\
\texttt{archive/\allowbreak ErdosStraus-\allowbreak core.lean} & 2073 & Frozen earlier snapshot of the covering file; not a Lake root. & \ref{subsec:legacy} \\
\texttt{archive/\allowbreak ErdosStraus-\allowbreak mathlib-\allowbreak project.lean} & 2480 & Duplicate of \file{ErdosStraus.lean}. & \ref{subsec:legacy} \\
\bottomrule
\end{longtable}
\endgroup
\noindent A short \texttt{archive/README.md} (not counted as Lean; see \Cref{subsec:markdown-docs})
accompanies the two frozen files and states explicitly that they are not Lake roots.

\subsection{Build configuration and licensing (9)}
\begin{itemize}[leftmargin=1.6em]
  \item \file{lakefile.toml} --- package \texttt{ErdosStraus}; declares the Mathlib dependency
        (pinned commit \texttt{1a0ef26d\dots}) and the eleven-file Layer-A/B library root list.
  \item \file{lean-toolchain} --- pins \texttt{leanprover/lean4:v4.34.0-rc1} (documented as a
        release candidate, not a moving nightly).
  \item \file{lake-manifest.json} --- resolved dependency manifest generated by \texttt{lake};
        \file{CONTRIBUTING.md} directs that it be committed with any toolchain or Mathlib bump.
  \item \file{.gitignore} --- excludes \texttt{.lake/}, \texttt{.oleans/}, \texttt{*.olean},
        \texttt{*.zip}, \texttt{\_\_pycache\_\_/}.
  \item \file{.github/workflows/ci.yml} --- GitHub Actions workflow; runs \texttt{lake exe cache
        get} and \texttt{lake build} against the pinned Mathlib cache on every push and pull
        request.
  \item \file{requirements.txt} --- a single pinned \texttt{numpy} version constraint for the
        numeric scripts of \Cref{sec:computation}.
  \item \file{CONTRIBUTING.md} --- pinning policy table, build instructions, and an explicit
        restatement of scope (no fake theorem, no \texttt{sorry} on the QED line, library files
        are not progress toward the load-bearing interface).
  \item \file{LICENSE} --- MIT licence, copyright Riley Betts Ltd, with an explicit carve-out
        noting that \lean{Stormer.lean} remains Apache~2.0.
  \item \file{LICENSE-APACHE} --- the Apache License~2.0 text covering \lean{Stormer.lean}
        alone, copyright Alexander Benjamin Worth Burns and Ravi Bajaj.
\end{itemize}

\subsection{Markdown documents (22)}
\label{subsec:markdown-docs}
\begingroup\small
\begin{longtable}{@{}p{0.30\linewidth}p{0.58\linewidth}p{0.06\linewidth}@{}}
\toprule
\textbf{File} & \textbf{Contents} & \textbf{\S} \\
\midrule
\endhead
\file{README.md} & Repository front matter: layer map, build instructions, axiom audit table,
  the ``novel results (unverified)'' table, the authorship/role statement, ``what's in Lean'' /
  ``what this does not claim'' summaries, full bibliography. & \ref{sec:lean} \\
\texttt{archive/\allowbreak README.md} & Short note: the two frozen Lean files are not Lake roots and
  should not be imported. & \ref{subsec:legacy} \\
\texttt{erdos-\allowbreak straus-\allowbreak programme.md} & The programme brief and repository front door: five open
  geometric questions, a chronological chain of dated postscripts (v0.9--v0.23) recording each
  investigation's outcome, the authorship statement. & \ref{subsec:whatitis},\ref{subsec:provenance} \\
\texttt{erdos-\allowbreak straus-\allowbreak novel-\allowbreak structures-\allowbreak plan.md} & The primary research log: eleven initial options
  considered, the programme (Tracks A/B/C), and sections \S4a--\S4w documenting the k-budget
  investigation, the Bright--Loughran route, and their synthesis. & \ref{subsec:kbudget},\ref{sec:route2} \\
\texttt{erdos-\allowbreak straus-\allowbreak road-\allowbreak to-\allowbreak lean-\allowbreak qed.md} & The execution roadmap: roadmap Layer A/B/C, the
  cardinality-to-QED bridge, the Renewal Lemma, and the target Lean architecture. & \ref{subsec:architecture} \\
\texttt{erdos-\allowbreak straus-\allowbreak conjecture-\allowbreak spec.md} & The acceptance specification, requirements R1--R8, and
  the scorecard of \Cref{tab:scorecard}. & \ref{sec:spec} \\
\texttt{erdos-\allowbreak straus-\allowbreak candidate-\allowbreak conjectures.md} & The full candidate list C1--C11, NV, S1, S1$_1$,
  with status and primary risk for each. & \ref{sec:candidates} \\
\file{erdos-straus-T-A.md} & T(A) working note: certificate architecture, defect layer, the two
  x-scan/growing-A measurements, the local Euler factor, T(A)$^+$, E\_lane. & \ref{subsec:t_a},\ref{subsec:talane} \\
\file{erdos-straus-T-3.md} & T(3) statement, proof of T(3)$^+$, the three-step lower-bound plan,
  and the $r_\chi$-to-Kloosterman look (negative result (iv)). & \ref{subsec:t3},\ref{subsec:t3-negative} \\
\file{erdos-straus-T-2.md} & The Zheng simultaneous-progressions look (negative result (i)). & \ref{subsec:t3-negative} \\
\texttt{erdos-\allowbreak straus-\allowbreak T-\allowbreak 3-\allowbreak delta.md} & The Heath-Brown $\delta$-method look (negative result (ii)). & \ref{subsec:t3-negative} \\
\texttt{erdos-\allowbreak straus-\allowbreak T-\allowbreak 3-\allowbreak monodromy.md} & The FKMS joint-monodromy look (negative result (iii)). & \ref{subsec:t3-negative} \\
\texttt{erdos-\allowbreak straus-\allowbreak varying-\allowbreak modulus-\allowbreak gap.md} & Capstone document: the four (with an informal
  fifth) T(3)-completion looks classified into one common failure mode. & \ref{subsec:varyingmodulus} \\
\texttt{zheng-\allowbreak wellfactorable.md} & Technical note on well-factorability of real characters (Zheng
  Def.~1.1), feeding \file{erdos-straus-T-2.md}. & \ref{subsec:t3-negative} \\
\texttt{erdos-\allowbreak straus-\allowbreak E-\allowbreak power.md} & Withdrawn one-stage write-up of Theorem~\ref{thm:epower};
  combinatorial core and recorded negative. See also \texttt{erdos-straus-E-power-decision.md}. & \ref{subsec:epower} \\
\texttt{erdos-\allowbreak straus-\allowbreak E-\allowbreak partial.md} & Gate~A write-up: what may/may not be claimed toward
  \lean{AnalyticSurvivorBound}; the schedule-retuning argument of \Cref{subsec:gatea}. & \ref{subsec:gatea} \\
\texttt{erdos-\allowbreak straus-\allowbreak gs-\allowbreak reformulation.md} & The Granville--Soundararajan restatement of the
  survivor-density target; the C2 null result. & \ref{subsec:gs} \\
\texttt{erdos-\allowbreak straus-\allowbreak sieve-\allowbreak desk.md} & The analytic question list. & \ref{subsec:sievedesk} \\
\texttt{erdos-\allowbreak straus-\allowbreak loughran-\allowbreak orbit.md} & The geometric question list: five numbered open questions
  in the Bright--Loughran literature. & \ref{subsec:loughran-questions} \\
\texttt{erdos-\allowbreak straus-\allowbreak prior-\allowbreak archive.md} & Triage of the sibling Track-1 covering-density archive: what
  was kept as library material and what was left out. & \ref{subsec:prior-archive} \\
\texttt{erdos-\allowbreak straus-\allowbreak bradford-\allowbreak lemma3.md} & Examination of the Bello-Hern\'andez--Benito--Fern\'andez
  rigidity lemma and of Bradford 2026 as a claimed proof. & \ref{subsec:bradford} \\
\texttt{erdos-\allowbreak straus-\allowbreak leochlon.md} & Documentation of the Leochlon Type-II identity dictionary. & \ref{subsec:leochlon} \\
\bottomrule
\end{longtable}
\endgroup

\subsection{Python scripts and their JSON outputs (20 scripts, 23 JSON files)}
\begingroup\small
\begin{longtable}{@{}p{0.26\linewidth}p{0.46\linewidth}p{0.20\linewidth}p{0.04\linewidth}@{}}
\toprule
\textbf{Script} & \textbf{Purpose} & \textbf{Output JSON(s)} & \textbf{\S} \\
\midrule
\endhead
\file{c4\_S\_xscan.py} & $x$-scan of $S(A,x)$ at fixed $A$; hard-prime sieve infrastructure
  reused throughout. & \file{c4\_S\_xscan\_X100000000.json}, \file{\_X1000000000.json} & \ref{sec:computation} \\
\file{c4\_growing\_A.py} & Growing-$A$ scan at fixed $x=10^9$; $\mathrm{cond}(a)$ and $\widehat C(A)$. &
  \file{c4\_growing\_A\_X1000000000\_A200.json} & \ref{sec:computation} \\
\file{c4\_growing\_box\_probe.py} & Earlier two-directional independence probe. &
  \file{c4\_growing\_box\_X1000000.json}, \file{\_X10000000.json} & \ref{sec:computation} \\
\file{c4\_two\_shift\_probe.py} & Pairwise-independence and C2 Jacobi-symbol probe. &
  \file{c4\_two\_shift\_probe.json} & \ref{sec:computation} \\
\file{c4\_euler\_factor.py} & Prime-aligned local Euler factor $C_{\mathrm{euler}}(A)$ and
  residue-collision check. & \file{c4\_euler\_factor\_Q10000000\_A200.json} & \ref{sec:computation} \\
\file{c4\_composite\_layer.py} & Genuine-extra composite Euler product $\Pi_{\mathrm{gen}}$. &
  \file{c4\_composite\_layer\_Z50000\_A80.json} & \ref{sec:computation} \\
\file{c4\_certificates.py} & Certificate converse and shared-certificate dummy-covering tests. &
  \file{c4\_certificates\_X1000000\_A40.json} & \ref{sec:computation} \\
\file{c4\_defect\_layer.py} & Exact/Monte-Carlo defect between \lean{ClassRough} and the
  certificate union. & \file{c4\_defect\_layer\_A12\_W5.json} & \ref{sec:computation} \\
\file{c4\_c2\_symbols.py} & C2 cross-moduli Jacobi-symbol correlation test. &
  \file{c4\_c2\_symbols\_X1000000\_A30.json}, \file{\_X100000000\_A80.json}, \file{\_X1000000000\_A80.json} & \ref{sec:computation} \\
\file{c4\_surface\_fit.py} & Joint $(A,x)$ surface fit of $\widehat C$; predictor comparison. &
  \file{c4\_surface\_fit.json} & \ref{sec:computation} \\
\file{c4\_sieve\_constant.py} & Selberg constant $C_{\mathrm{sieve}}(A)$ and $B(A)$ of T(A)$^+$. &
  \file{c4\_sieve\_constant\_Q1000000\_A200.json}, \file{\_Q10000000\_A200.json} & \ref{sec:computation} \\
\file{c4\_t3\_euler.py} & T(3)-specific certificate-assignment Euler products and occupancy. &
  \file{c4\_t3\_euler\_Q1000000.json} & \ref{sec:computation} \\
\file{e\_power\_suen\_moments.py} & Lemma~SM ratio checks, $A\le240$. &
  \file{e\_power\_suen\_moments\_A240.json} & \ref{subsec:computation-epower} \\
\file{e\_power\_suen\_moments\_large.py} & Lemma~SM ratio checks extended to $A\le2000$. &
  \file{e\_power\_suen\_moments\_A2000.json} & \ref{subsec:computation-epower} \\
\file{t2\_zheng\_ranges.py} & Zheng $\mathcal L(\theta)$ range check (negative result (i)). &
  \file{t2\_zheng\_ranges.json} & \ref{subsec:t3-negative} \\
\file{t3\_delta\_ranges.py} & Farey--$k$ incompatibility check (negative result (ii)). &
  \file{t3\_delta\_ranges.json} & \ref{subsec:t3-negative} \\
\file{t3\_kloosterman\_ranges.py} & Pascadi/Blomer--Pascadi/MQW saving-window check (negative result (iv)). &
  \file{t3\_kloosterman\_ranges.json} & \ref{subsec:t3-negative} \\
\file{t3\_monodromy\_ranges.py} & FKMS gallant-monodromy length check (negative result (iii)). &
  \file{t3\_monodromy\_ranges.json} & \ref{subsec:t3-negative} \\
\file{s1\_schinzel\_separators.py} & Schinzel narrow-band separator sweep. & (no separate JSON; prints results) & \ref{subsec:schinzel} \\
\file{c7\_bounded\_coupling\_search.py} & C7 identity search and coupling-mass comparison. & (no separate JSON; prints results) & \ref{subsec:computation-misc} \\
\bottomrule
\end{longtable}
\endgroup
\noindent The count of 23 JSON files above excludes \file{lake-manifest.json}
(\Cref{app:files}.2), which is a build artefact rather than an experimental output.


\section{The axiom audit}
\label{app:axioms}

\Cref{tab:axioms} reproduces, for reference, the axiom audit table maintained in
\file{README.md}, obtained by \texttt{\#print axioms} on Lean~4.33.0 for every named theorem on
the file-map Layer-A kernel-checked spine. \texttt{native\_decide} instances add a generated
reduction axiom (\texttt{Lean.ofReduceBool}, materialising as
\texttt{\_.\_native.native\_decide.ax\_1\_1} on 4.33.0) and are listed separately.

\begin{table}[h]
\centering\small
\begin{tabular}{@{}ll@{}}
\toprule
Theorem & Axioms \\
\midrule
\lean{witness\_sound} & \texttt{propext}, \texttt{Quot.sound} \\
\lean{es\_1009} & \texttt{propext}, \texttt{Quot.sound} \\
\lean{bl\_thm12\_prime\_anatomy} & \texttt{propext}, \texttt{Quot.sound} \\
\lean{primes\_below\_100\_covered\_kernel} & none (\texttt{decide}) \\
\lean{lemma38\_check\_true} & none (\texttt{decide}) \\
\lean{conditional\_qed\_hard} & \texttt{propext}, \texttt{Quot.sound}, \texttt{Classical.choice} \\
\lean{lemma32} & \texttt{propext}, \texttt{Quot.sound}, \texttt{Classical.choice} \\
\lean{ConicFiber.fiber\_identity} & \texttt{propext}, \texttt{Quot.sound} \\
\lean{ConicFiber.fiber\_to\_divisor} & \texttt{propext}, \texttt{Quot.sound} \\
\lean{ConicFiber.divisor\_to\_fiber} & \texttt{propext}, \texttt{Quot.sound} \\
\lean{TwistDescent.descent} & none \\
\lean{TwistDescent.naive\_cover\_empty} & \texttt{propext}, \texttt{Quot.sound} \\
\lean{TwistDescent.lift} & \texttt{propext}, \texttt{Quot.sound}, \texttt{Classical.choice} \\
\lean{TwistDescent.twist\_real\_iff} & \texttt{propext} \\
\lean{TwistDescent.t\_fiber\_section\_cleared} & \texttt{propext} \\
\lean{TwistDescent.u1\_fiber\_section\_cleared} & \texttt{propext} \\
\lean{TwistDescent.u1\_fiber\_twist} & \texttt{propext} \\
\lean{SchinzelSep.signed\_9\_5} & none \\
\lean{SchinzelSep.no\_pos\_9\_5} & \texttt{propext}, \texttt{Quot.sound} \\
\lean{SchinzelDecide.sep\_9\_5} & none (\texttt{decide}) \\
\lean{SchinzelDecide.sep\_13\_7} & none (\texttt{decide}) \\
\lean{SchinzelDecide.pos\_4\_5} & none (\texttt{decide}) \\
\lean{SchinzelDecide.decideCplus\_sound} & \texttt{propext}, \texttt{Quot.sound} \\
\lean{SchinzelDecide.decideCplus\_complete} & \texttt{propext}, \texttt{Quot.sound} \\
\lean{SchinzelDecide.no\_pos\_9\_5} & \texttt{propext}, \texttt{Classical.choice}, \texttt{Quot.sound} \\
\lean{SchinzelDecide.no\_pos\_13\_7} & \texttt{propext}, \texttt{Classical.choice}, \texttt{Quot.sound} \\
\lean{NoVieta.unique\_x} & \texttt{propext}, \texttt{Quot.sound} \\
\lean{NoVieta.es\_unique\_x} & \texttt{propext}, \texttt{Classical.choice}, \texttt{Quot.sound} \\
\bottomrule
\end{tabular}
\caption{Axiom audit of the kernel-checked spine (\file{README.md}).}
\label{tab:axioms}
\end{table}

\noindent The five \texttt{native\_decide} theorems, each additionally carrying
\texttt{Lean.ofReduceBool}, are: \lean{primes\_below\_10000\_covered},
\lean{hard\_primes\_below\_10000\_covered} and \lean{okay\_below\_100000} in
\file{ErdosStraus.lean}, and \lean{thm12\_instance\_1009},
\lean{witness\_fingerprint\_1009} in \file{BrightLoughran.lean}.


\addcontentsline{toc}{section}{References}
\begin{thebibliography}{99}
\small

\bibitem{Erdos1950} P.~Erd\H{o}s, \emph{Az $1/x_1+\cdots+1/x_n=a/b$ egyenlet eg\'esz sz\'am\'u
megold\'asair\'ol}, Mat.\ Lapok \textbf{1} (1950), 192--210.

\bibitem{Sierpinski1956} W.~Sierpi\'nski, \emph{Sur les d\'ecompositions de nombres rationnels en
fractions primaires}, Mathesis \textbf{65} (1956), 16ff.

\bibitem{Schinzel1956} A.~Schinzel, \emph{Sur quelques propri\'et\'es des nombres $3/n$ et $4/n$,
o\`u $n$ est un nombre impair}, Mathesis \textbf{65} (1956), 219--222.

\bibitem{Nakayama1939} N.~Nakayama, \emph{On the decomposition of a rational number into
``Stammbr\"uche''}, T\^ohoku Math.\ J. \textbf{46} (1939), 1ff.

\bibitem{Rosati1954} L.~A.~Rosati, \emph{Sull'equazione diofantea $4/n=1/x_1+1/x_2+1/x_3$},
Boll.\ Un.\ Mat.\ Ital.\ (3) \textbf{9} (1954), 59ff.

\bibitem{Aigner1964} A.~Aigner, \emph{Br\"uche als Summe von Stammbr\"uchen}, J.\ Reine Angew.\
Math.\ \textbf{214/215} (1964), 174ff.

\bibitem{Yamamoto1965} K.~Yamamoto, \emph{On the Diophantine equation $4/n=1/x+1/y+1/z$}, Mem.\
Fac.\ Sci.\ Kyushu Univ.\ Ser.\ A \textbf{19} (1965), 37--47.

\bibitem{Mordell1969} L.~J.~Mordell, \emph{Diophantine Equations}, Academic Press, London, 1969,
Chapter~30.

\bibitem{Schinzel2000} A.~Schinzel, \emph{On sums of three unit fractions with polynomial
denominators}, Funct.\ Approx.\ Comment.\ Math.\ \textbf{28} (2000), 187--194.

\bibitem{Bloom2022} T.~F.~Bloom, \emph{Egyptian fractions}, arXiv:2210.04496.

\bibitem{Lopez2024} M.~A.~Lopez, \emph{A complete congruence system for the Erd\H{o}s--Straus
conjecture}, arXiv:2404.01508.

\bibitem{BelloHernandez2026} M.~Bello-Hern\'andez, M.~Benito, E.~Fern\'andez, \emph{A divisor
parametrization for the Erd\H{o}s--Straus conjecture}, arXiv:2606.10922.

\bigskip\noindent\textbf{Bright--Loughran geometry (Route 2).}

\bibitem{BrightLoughran2020} M.~Bright, D.~Loughran, \emph{Brauer--Manin obstruction for
Erd\H{o}s--Straus surfaces}, Bull.\ Lond.\ Math.\ Soc.\ \textbf{52} (2020), no.~4, 746--761;
arXiv:1908.02526.

\bibitem{Serre1973} J.-P.~Serre, \emph{Cours d'arithm\'etique}, PUF, Paris, 1970; English:
\emph{A Course in Arithmetic}, GTM~7, Springer, 1973.

\bibitem{LangWeil1954} S.~Lang, A.~Weil, \emph{Number of points of varieties in finite fields},
Amer.\ J.\ Math.\ \textbf{76} (1954), 819--827.

\bibitem{JahnelSchindler2017} J.~Jahnel, D.~Schindler, \emph{On integral points on degree four
del Pezzo surfaces}, Israel J.\ Math.\ \textbf{222} (2017), 21--62; arXiv:1602.03118.

\bibitem{Harpaz2019} Y.~Harpaz, \emph{Integral points on conic log K3 surfaces}, J.\ Eur.\ Math.\
Soc.\ \textbf{21} (2019), 627--664; arXiv:1511.04876.

\bibitem{Uppal2023} H.~Uppal, \emph{Integral points on affine surfaces fibered over $\mathbb A^1$},
arXiv:2307.06638.

\bibitem{CTWittenberg2012} J.-L.~Colliot-Th\'el\`ene, O.~Wittenberg, \emph{Groupe de Brauer et
points entiers de deux familles de surfaces cubiques affines}, Amer.\ J.\ Math.\ \textbf{134}
(2012), 1303--1327.

\bibitem{CaoXu2018} Y.~Cao, F.~Xu, \emph{Strong approximation with Brauer--Manin obstruction for
toric varieties}, Ann.\ Inst.\ Fourier \textbf{68} (2018), 1879--1908.

\bibitem{CTXu2009} J.-L.~Colliot-Th\'el\`ene, F.~Xu, \emph{Brauer--Manin obstruction for integral
points of homogeneous spaces and representation by integral quadratic forms}, Compos.\ Math.\
\textbf{145} (2009), 309--363.

\bigskip\noindent\textbf{Markoff comparison.}

\bibitem{BGS2016} J.~Bourgain, A.~Gamburd, P.~Sarnak, \emph{Markoff triples and strong
approximation}, C.\ R.\ Math.\ Acad.\ Sci.\ Paris \textbf{354} (2016), 131--135;
arXiv:1505.06411; sequel \emph{Markoff surfaces and strong approximation: 1},
arXiv:1607.01530.

\bibitem{GhoshSarnak2022} A.~Ghosh, P.~Sarnak, \emph{Integral points on Markoff type cubic
surfaces}, Invent.\ Math.\ \textbf{229} (2022), 689--749; arXiv:1706.06712.

\bibitem{LoughranMitankin2021} D.~Loughran, V.~Mitankin, \emph{Integral Hasse principle and
strong approximation for Markoff surfaces}, Int.\ Math.\ Res.\ Not.\ IMRN \textbf{2021}, no.~18,
14086--14122; arXiv:1807.10223.

\bibitem{CTWX2020} J.-L.~Colliot-Th\'el\`ene, D.~Wei, F.~Xu, \emph{Brauer--Manin obstruction for
Markoff surfaces}, Ann.\ Sc.\ Norm.\ Super.\ Pisa Cl.\ Sci.\ (5) \textbf{21} (2020), 1257--1313;
arXiv:1808.01584.

\bibitem{MitankinUhlemann2025} V.~Mitankin, J.~Uhlemann, \emph{Local--global principles for
semi-integral points on Markoff orbifold pairs}, J.\ Lond.\ Math.\ Soc.\ \textbf{112} (2025),
no.~6, e70363; arXiv:2411.02629.

\bibitem{Jang2023} S.~Jang, \emph{Tropical dynamics of Markov surfaces}, arXiv:2306.11357.

\bigskip\noindent\textbf{Analytic exceptional set and sieve theory.}

\bibitem{Vaughan1970} R.~C.~Vaughan, \emph{On a problem of Erd\H{o}s, Straus and Schinzel},
Mathematika \textbf{17} (1970), 193--198.

\bibitem{Suen1990} W.-C.~S.~Suen, \emph{A correlation inequality and a Poisson limit theorem for
nonoverlapping balanced subgraphs of a random graph}, Random Structures Algorithms \textbf{1}
(1990), 231--242.

\bibitem{Janson1990} S.~Janson, \emph{Poisson approximation for large deviations}, Random
Structures Algorithms \textbf{1} (1990), 221--229.

\bibitem{ElsholtzTao2013} C.~Elsholtz, T.~Tao, \emph{Counting the number of solutions to the
Erd\H{o}s--Straus equation on unit fractions}, J.\ Aust.\ Math.\ Soc.\ \textbf{94} (2013), 50--105.

\bibitem{TaoBlog2011} T.~Tao, \emph{On the number of solutions to $4/p=1/n_1+1/n_2+1/n_3$}, blog
post, 19 November 2011.

\bibitem{Shiu1980} P.~Shiu, \emph{A Brun--Titchmarsh theorem for multiplicative functions}, J.\
Reine Angew.\ Math.\ \textbf{313} (1980), 161--170.

\bibitem{Nair1992} M.~Nair, \emph{Multiplicative functions of polynomial values in short
intervals}, Acta Arith.\ \textbf{62} (1992), 257--269.

\bibitem{NairTenenbaum1998} M.~Nair, G.~Tenenbaum, \emph{Short sums of certain arithmetic
functions}, Acta Math.\ \textbf{180} (1998), 119--144.

\bibitem{Henriot2012} K.~Henriot, \emph{Nair--Tenenbaum bounds uniform with respect to the
discriminant}, Math.\ Proc.\ Cambridge Philos.\ Soc.\ \textbf{152} (2012), 405--424; erratum
\textbf{157} (2014); arXiv:1102.1643.

\bibitem{FIMR2013} J.~B.~Friedlander, H.~Iwaniec, B.~Mazur, K.~Rubin, \emph{The spin of prime
ideals}, Invent.\ Math.\ \textbf{193} (2013), 697--749.

\bibitem{Maire2018} C.~Maire, \emph{Genus theory and governing fields}, New York J.\ Math.\
\textbf{24} (2018), 1056--1067.

\bibitem{HalberstamRichert1974} H.~Halberstam, H.-E.~Richert, \emph{Sieve Methods}, Academic
Press, London, 1974.

\bibitem{OperaDeCribro} J.~Friedlander, H.~Iwaniec, \emph{Opera de Cribro}, AMS Colloquium
Publications \textbf{57}, 2010.

\bibitem{Norton1976} K.~K.~Norton, \emph{On the number of restricted prime factors of an integer.
I}, Illinois J.\ Math.\ \textbf{20} (1976), 681--705.

\bibitem{LanguascoZaccagnini2007} A.~Languasco, A.~Zaccagnini, \emph{A note on Mertens' formula
for arithmetic progressions}, J.\ Number Theory \textbf{127} (2007), 37--46.

\bibitem{GranvilleSoundararajan2003} A.~Granville, K.~Soundararajan, \emph{The distribution of
values of $L(1,\chi_d)$}, Geom.\ Funct.\ Anal.\ \textbf{13} (2003), 992--1028.

\bibitem{GranvilleSoundararajan2001} A.~Granville, K.~Soundararajan, \emph{Large character sums},
J.\ Amer.\ Math.\ Soc.\ \textbf{14} (2001), 365--397.

\bibitem{GranvilleSoundararajan2007} A.~Granville, K.~Soundararajan, \emph{Pretentious
multiplicative functions and an inequality for the zeta-function}, 2007, and subsequent
pretentiousness papers.

\bibitem{HeathBrown1995} D.~R.~Heath-Brown, \emph{A mean value estimate for real character sums},
Acta Arith.\ \textbf{72} (1995), 235--275.

\bibitem{Iwaniec1974} H.~Iwaniec, \emph{Primes represented by quadratic polynomials in two
variables}, Acta Arith.\ \textbf{24} (1973/74), 435--459.

\bibitem{Iwaniec1976} H.~Iwaniec, \emph{The half dimensional sieve}, Acta Arith.\ \textbf{29}
(1976), 69--95.

\bibitem{BrudernFouvry1994} J.~Br\"udern, \'E.~Fouvry, \emph{Lagrange's four squares theorem with
almost prime variables}, J.\ Reine Angew.\ Math.\ \textbf{454} (1994), 59--96.

\bibitem{BrudernFouvry1996} J.~Br\"udern, \'E.~Fouvry, \emph{Le crible \`a vecteurs}, Compositio
Math.\ \textbf{102} (1996), 337--355.

\bibitem{Harman2007} G.~Harman, \emph{Prime-detecting sieves}, LMS Monographs \textbf{33},
Princeton, 2007.

\bibitem{Linnik1960} Ju.~V.~Linnik, \emph{An asymptotic formula in an additive problem of
Hardy--Littlewood}, Izv.\ Akad.\ Nauk SSSR Ser.\ Mat.\ \textbf{24} (1960), 629--706.

\bibitem{Hooley1973} C.~Hooley, \emph{On the intervals between numbers that are sums of two
squares: II}, J.\ Number Theory \textbf{5} (1973), 215--217.

\bibitem{Zheng2025} Z.~Zheng, \emph{Primes in simultaneous arithmetic progressions},
arXiv:2512.22798.

\bibitem{BombieriFriedlanderIwaniec1986} E.~Bombieri, J.~B.~Friedlander, H.~Iwaniec, \emph{Primes
in arithmetic progressions to large moduli}, Acta Math.\ \textbf{156} (1986), 203--251.

\bibitem{HeathBrown1978} D.~R.~Heath-Brown, \emph{Hybrid bounds for Dirichlet $L$-functions},
Invent.\ Math.\ \textbf{47} (1978), 149--170.

\bibitem{GrimmeltMerikoski2024} L.~Grimmelt, J.~Merikoski, \emph{Twisted correlations of the
divisor function via discrete averages of $\mathrm{SL}_2(\mathbb R)$ Poincar\'e series},
arXiv:2404.08502.

\bibitem{DeshouillersIwaniec1982} J.-M.~Deshouillers, H.~Iwaniec, \emph{Kloosterman sums and
Fourier coefficients of cusp forms}, Invent.\ Math.\ \textbf{70} (1982/83), 219--288.

\bibitem{FouvryIwaniec1997} \'E.~Fouvry, H.~Iwaniec, \emph{Gaussian primes}, Acta Arith.\
\textbf{79} (1997), 249--287.

\bibitem{KowalskiMichelSawin2017} E.~Kowalski, P.~Michel, W.~Sawin, \emph{Bilinear forms with
Kloosterman sums and applications}, Ann.\ of Math.\ \textbf{186} (2017), 413--500.

\bibitem{BlomerMilicevic2015} V.~Blomer, D.~Mili\'cevi\'c, \emph{The second moment of twisted
modular $L$-functions}, Geom.\ Funct.\ Anal.\ \textbf{25} (2015), 453--516.

\bibitem{Shkredov2021} I.~D.~Shkredov, \emph{Modular hyperbolas and bilinear forms of Kloosterman
sums}, J.\ Number Theory \textbf{220} (2021), 182--211.

\bibitem{Pascadi2026} A.~Pascadi, \emph{Non-abelian amplification and bilinear forms with
Kloosterman sums}, Geom.\ Funct.\ Anal., accepted; arXiv:2511.08445.

\bibitem{MQW2026} D.~Mili\'cevi\'c, Y.~Qin, H.~Wu, \emph{Bilinear forms with Kloosterman sums and
moments of twisted $L$-functions}, arXiv:2511.07550.

\bibitem{BlomerPascadi2026} V.~Blomer, A.~Pascadi, \emph{Bilinear forms with Kloosterman sums via
quadratic characters}, arXiv:2607.24311.

\bibitem{NathXie2025} K.~Nath, L.~Xie, \emph{Almost primes and primes that are sums of two
squares plus one}, arXiv:2501.16723; Acta Arith.\ 2025.

\bibitem{FHRSS2025} E.~Fuchs, C.~Hsu, J.~Rickards, D.~Schindler, K.~E.~Stange, \emph{Primes
represented by shifted quadratic forms: on primitivity and congruence classes},
arXiv:2504.20289; Acta Arith.\ \textbf{222} (2026), 371--391.

\bigskip\noindent\textbf{The $\delta$-symbol, two-dimensional Kloosterman sums, and joint
monodromy (T(3)-completion investigations).}

\bibitem{HeathBrown1996} D.~R.~Heath-Brown, \emph{A new form of the circle method, and its
application to quadratic forms}, J.\ Reine Angew.\ Math.\ \textbf{481} (1996), 149--206.

\bibitem{DFI} W.~Duke, J.~B.~Friedlander, H.~Iwaniec, \emph{Equidistribution of roots of a
quadratic congruence to prime moduli}, Ann.\ of Math.\ \textbf{141} (1995), 423--441; and
\emph{Class group L-functions}, Invent.\ Math.\ \textbf{128} (1997), 23--43 (used in
\texttt{erdos-\allowbreak straus-\allowbreak T-\allowbreak 3-\allowbreak delta.md} for the two-dimensional analogue).

\bibitem{HeathBrownPierce2017} D.~R.~Heath-Brown, L.~B.~Pierce, \emph{Simultaneous integer values
of pairs of quadratic forms}, J.\ Reine Angew.\ Math.\ \textbf{727} (2017), 85--143;
arXiv:1309.6767.

\bibitem{BettinChandee2018} S.~Bettin, V.~Chandee, \emph{Trilinear forms with Kloosterman
fractions}, Adv.\ Math.\ \textbf{328} (2018), 1234--1262.

\bibitem{Katz1988} N.~M.~Katz, \emph{Gauss Sums, Kloosterman Sums, and Monodromy Groups}, Ann.\
of Math.\ Studies \textbf{116}, Princeton University Press, 1988.

\bibitem{FKMS2026} \'E.~Fouvry, E.~Kowalski, Ph.~Michel, W.~Sawin, \emph{Bilinear forms with
trace functions}, arXiv:2511.09459 (v3, 11 August 2026).

\bibitem{FouvryKowalskiMichel2014} \'E.~Fouvry, E.~Kowalski, Ph.~Michel, \emph{Algebraic trace
functions over the primes}, Duke Math.\ J.\ \textbf{163} (2014), 1683--1736.

\bibitem{FouvryKowalskiMichel2015} \'E.~Fouvry, E.~Kowalski, Ph.~Michel, \emph{A study in sums
of products}, Philos.\ Trans.\ Roy.\ Soc.\ A \textbf{373} (2015), 20140309.

\bibitem{RicottaRoyer2016} G.~Ricotta, E.~Royer, \emph{Kloosterman paths of prime powers
moduli}, arXiv:1609.03694.

\bibitem{Kowalski2006} E.~Kowalski, \emph{The large sieve, monodromy, and zeta functions of
curves}, arXiv:math/0608718 (big symplectic/orthogonal monodromy modulo $\ell$, as both the
field and a conductor grow).

\bigskip\noindent\textbf{Prior archive and external claims examined (\Cref{sec:prior}).}

\bibitem{Stormer1897} C.~St\o rmer, \emph{Quelques th\'eor\`emes sur l'\'equation de Pell
$x^2-Dy^2=\pm1$ et leurs applications}, Videnskabsselskabets Skrifter I Math.-Nat.\ Kl.\
\textbf{2} (1897).

\bibitem{BelloHernandez2010} M.~Bello-Hern\'andez, M.~Benito, E.~Fern\'andez, \emph{On the
equation $4/p=1/n+1/m+1/k$}, arXiv:1010.2035.

\bibitem{Bradford2026} K.~Bradford, \emph{An elementary proof of the Erd\H{o}s--Straus
conjecture}, arXiv:2602.11774.

\bibitem{Bradford2021} K.~Bradford, \emph{Integers} \textbf{21} (2021), Paper No.~A24.

\bibitem{BradfordIonascu2015} K.~Bradford, E.~J.~Ionascu, 2015 (as cited in
\texttt{erdos-\allowbreak straus-\allowbreak bradford-\allowbreak lemma3.md}; full bibliographic
details not reproduced from the source document).

\bibitem{Bradford2024} K.~Bradford, arXiv:2403.16047.

\bigskip\noindent\textbf{Computation.}

\bibitem{Salez2014} S.~E.~Salez, \emph{The Erd\H{o}s--Straus conjecture: new modular equations
and checking up to $N=10^{17}$}, arXiv:1406.6307.

\end{thebibliography}

\end{document}
